%% LyX 2.4.3 created this file.  For more info, see https://www.lyx.org/.
%% Do not edit unless you really know what you are doing.
\documentclass[11pt,english]{article}
\usepackage[T1]{fontenc}
\usepackage[latin9]{inputenc}
\usepackage{xcolor}
\usepackage{float}
\usepackage{amsmath}
\usepackage{amssymb}
\usepackage{geometry}
\geometry{verbose,tmargin=1in,bmargin=1in,lmargin=1in,rmargin=1in}
\usepackage{setspace}
\usepackage[authoryear]{natbib}
\usepackage[all]{xy}
\PassOptionsToPackage{normalem}{ulem}
\usepackage{ulem}
\onehalfspacing

\makeatletter

%%%%%%%%%%%%%%%%%%%%%%%%%%%%%% LyX specific LaTeX commands.
\providecolor{lyxadded}{rgb}{0,0,1}
\providecolor{lyxdeleted}{rgb}{1,0,0}
%% Change tracking with ulem and xcolor: base macros
\DeclareRobustCommand{\mklyxadded}[1]{\textcolor{lyxadded}\bgroup#1\egroup}
\DeclareRobustCommand{\mklyxdeleted}[1]{\textcolor{lyxdeleted}\bgroup\mklyxsout{#1}\egroup}
\DeclareRobustCommand{\mklyxsout}[1]{\ifx\\#1\else\sout{#1}\fi}
%% Change tracking with ulem and xcolor: ct markup
\DeclareRobustCommand{\lyxadded}[4][]{\mklyxadded{#4}}
\DeclareRobustCommand{\lyxdeleted}[4][]{\mklyxdeleted{#4}}

%%%%%%%%%%%%%%%%%%%%%%%%%%%%%% User specified LaTeX commands.
\usepackage{bbm}
\usepackage{float}
\usepackage{siunitx}
\usepackage{bbm}
\usepackage{amsfonts}
\usepackage{amsmath}
\usepackage{graphicx}
\usepackage{setspace}
\usepackage[authoryear]{natbib}
\usepackage{subcaption}
\usepackage{booktabs}
\usepackage{color}
\usepackage[maxfloats=36]{morefloats}
\usepackage{geometry}
\usepackage{pdflscape}
\usepackage{tabularx}
\usepackage{multirow}
% Formatting
\usepackage{charter}
\usepackage[T1]{fontenc}
\geometry{verbose,tmargin=1.125in,bmargin=1.125in,lmargin=1.125in,rmargin=1.125in}


\date{}
\usepackage{datetime}
\newdateformat{monthyeardate}{%
\monthname[\THEMONTH] \THEYEAR
}
\newcommand{\ssize}{5.5cm}
\newcommand{\mssize}{6cm}
\newcommand{\msize}{7.5cm}
\newcommand{\lsize}{9cm}
\newcommand{\vlsize}{11cm}

\usepackage{rotating}

% Some math help
\DeclareMathOperator*{\argmax}{argmax}
\newcommand{\E}[0]{\mathbb{E}}

% For inputting scalars
\usepackage{xstring}
\usepackage{ifthen}
\usepackage{catchfile}

\newcommand{\scalarinput}[1]{\input{../output/tables/scalars/#1}\unskip} % Keep minus signs
\newcommand{\absinput}[1]{%
    \CatchFileDef{\filecontents}{../output/tables/scalars/#1}{}% Read the file contents into \filecontents
    \IfBeginWith{\filecontents}{-}{% Check if the first character is a minus sign
        \StrGobbleLeft{\filecontents}{1}[\filecontents]% Remove the first character
    }{}%
    \filecontents\unskip% Output the processed contents
}
%\usepackage{titlesec}
%\titlespacing*{\section} {0pt}{2.9ex plus 1ex minus .2ex}{1.8ex plus .2ex}
%\titlespacing*{\subsection} {0pt}{2.4ex plus 1ex minus .2ex}{1.5ex plus .2ex}

\usepackage{titling}
\setlength{\droptitle}{-2cm}   % adjust -2cm up or down to taste

\newcommand{\sizecutoff}{\$100,000 }


\usepackage{authblk}
\usepackage{lipsum} % for dummy text

% Reduce space between author blocks
\renewcommand\Authsep{\hspace{4em}}
\renewcommand\Authands{\hspace{4em}}
\renewcommand\Authfont{\normalsize}
\renewcommand\Affilfont{\itshape\small}
\usepackage{makecell}

\makeatother

\usepackage{babel}
\begin{document}
\renewcommand*{\thefootnote}{\fnsymbol{footnote}}
\title{Will that Be Cash or Credit: A Quantitative Analysis of Modern Interchange
Fees \thanks{We thank seminar and conference participants at Emory University,
Northwestern University, and Stanford SITE. We thank Fiserv for graciously
providing the data. Vincent Yao is an unpaid member of the Fiserv
Small Business Index Advisory Council and a paid consultant. Mark
Egan is an unpaid member of the Fiserv Small Business Index Advisory
Council.}}
\author{
\makebox[\textwidth][c]{%
  \begin{minipage}{0.33\textwidth}\centering
    Mark Egan\thanks{Harvard University, Harvard Business School. Email: megan@hbs.edu.} \\ \vspace{-1em}
    Harvard University and NBER
  \end{minipage}
  \hfill
  \begin{minipage}{0.33\textwidth}\centering
    Gregor Matvos\thanks{Northwestern University, Kellogg School of Management. Email: gregor.matvos@kellogg.northwestern.edu.} \\ \vspace{-1em}
    Northwestern University and NBER
  \end{minipage}
  \hfill
  \begin{minipage}{0.33\textwidth}\centering
    Amit Seru\thanks{Stanford GSB. Email: aseru@stanford.edu.} \\ \vspace{-1em}
    Stanford University and NBER
  \end{minipage}
}

\vspace{2em}

% Second row of authors
\makebox[\textwidth][c]{%
  \begin{minipage}{0.33\textwidth}\centering
    Lulu Wang\thanks{Northwestern University, Kellogg School of Management. Email: lulu.wang@kellogg.northwestern.edu.} \\ \vspace{-1em}
    Northwestern University
  \end{minipage}
  
  \begin{minipage}{0.33\textwidth}\centering
    Vincent Yao\thanks{Georgia State University, J. Mack Robinson College of Business. Email: wyao2@gsu.edu.} \\ \vspace{-1em}
    Georgia State University
  \end{minipage}
}}

\maketitle
\begin{center}
\vspace*{-1.5cm}
\textcolor{red}{PRELIMINARY AND INCOMPLETE. PLEASE DO NOT SHARE!}
\end{center} 
\pagenumbering{roman} \thispagestyle{empty} \vspace{0.5cm} 
    
\begin{abstract}
\begin{singlespace}
We examine how the fees that merchants pay to process card payments
shape retail prices and redistribute consumption. Credit cards, especially
premium ones, offer generous rewards but impose higher interchange
fees on merchants.�If merchants pass through these fees to retail
prices, users of low-cost payment methods (e.g., cash, debit) effectively
cross-subsidize high-reward credit card users. Using large-scale merchant-level
data that covers both card and cash transactions, we document substantial
heterogeneity in payment composition and interchange fees across merchants,
as well as a causal impact of fees on prices. �Variation in average
interchange fees leads to differences in retail prices, suggesting
that low-cost, low-reward users subsidize high-reward cardholders.
However, consumer sorting---where shoppers who use different payment
methods also shop at different merchants---mitigates this cross-subsidization.�Additionally,
merchants serving a broad range of payment types, such as grocery
stores, negotiate the lowest interchange rates, further reducing redistribution.
We embed these results in a sufficient-statistics framework to quantify
how the payment system redistributes consumption via prices and card
rewards. We estimate that interchange fees generate an implicit annual
transfer of approximately \$\scalarinput{combined_dollar_headline}
billion from cash and debit users to credit card users --- comparable
in scale but opposite in direction to major government transfer programs.
�Consumer sorting and negotiated rates reduce regressivity substantially
by \scalarinput{pct_reduction}\%, but do not eliminate it. Our findings
highlight how the rise of premium credit cards has reshaped the redistribution
of consumption, with significant implications for competition, regulation,
and the distributional impact of financial innovation.
\end{singlespace}

\thispagestyle{empty} \setcounter{page}{0}\begin{flushleft}

Keywords: Interchange Fees, Redistribution, Durbin Amendment, Price
Discrimination, Payment Systems

JEL Classification: E42, D14, L11, L81.\end{flushleft}
\end{abstract}
\clearpage{}

\pagenumbering{arabic} \renewcommand*{\thefootnote}{\arabic{footnote}}
\setcounter{footnote}{0}

We need to say much more about where the re-distributoin is largest:

- big vs. small

- basic goods vs luxory goods

- cheap or expensive stores

- etc. etc.

Link to important topcis in liteatrue:

- how does shopping / search behavior affect these transfers?-->
if cheaper stores

- what has happened to reddistribution over time
\begin{itemize}
\item Where is the relative contribution to the literature that shrouding
etc. results in redistribution
\begin{itemize}
\item normally the idea in this liteature is that everyone faces one pooled
contract (think insurance, refinance etc) and that this pooling results
in MORE redistribution; in our setting, pooling people on same payment
method decreases redistribution
\begin{itemize}
\item pooling on same payment is different than pooling on same merchant
\end{itemize}
\item sorting, negotationed fees etc. matters because it modulates this
redistribution
\begin{itemize}
\item what is new / surprising about this?
\item does it change how we think about policy?
\end{itemize}
\end{itemize}
\end{itemize}

\section{Introduction}

Payment systems can result in transfers across different groups of
consumers who choose different payment methods. Often termed \textquotedblleft reverse
Robin Hood\textquotedblright{} redistribution arises because merchant
fees for accepting payments vary significantly across payment types
yet retail prices typically do not differ by payment method \citep{Stavins2018}.
In turn, the major component of merchant fees--the interchange fee--is
funneled back to consumers though rewards. As merchants prices are
partially set to cover the cost of payments, consumers who use no
or low reward payment methods the rewards enjoyed by some credit card
users \citep{carlton1994antitrust}. These redistributive effects
have drawn increasing scrutiny from policymakers, most recently with
the introduction of the Credit Card Competition Act \citep{Andriotis2022}
\footnote{Some other relevant regulations include the Durbin Act of 2011, the
North Dakota federal court vacating the ruling in 2025 which is pending
appeal, or the Colorado\textquoteright s \textquotedblleft Swipe Fee
Fairness and Consumer Safeguards Act\textquotedblright{} at the state
level} The U.S. strongly contrasts with the European Union\textquoteright s
interchange fee regulation (IFR), which caped interchange fees less
than 10\% of the highest fees charged in the U.S.. Several features
of the industry have amplified these concerns. One is the growth of
premium credit cards, such as the Chase Sapphire Reserve, which are
funded by higher interchange fees and provide higher rewards. The
second is the ability of large merchants to negotiate interchange
fees, while small merchants are charged the standard, (higher) rates
following a set schedule. 

Despite their relevance, the effects of interchange fees on retail
pricing and redistribution have been challenging to study due to a
lack of merchant-level data on payment composition and acceptance
costs. These data are essential because the welfare consequences of
interchange fees depend not only on the magnitude of fee pass-through
to prices, but also on where consumers shop \citep{Wang.etal2014,Gans2018}.
Redistribution arises only when consumers using low-cost payment methods
(e.g., cash or debit) shop at the same merchants as those using high-cost
methods (e.g., rewards credit cards). Thus, the total fees a merchant
pays---and the extent of redistribution, which occurs at the merchant
level---depend on the specific mix of payment methods and their associated
acceptance costs. While card networks publish headline interchange
fee schedules, the actual fees incurred are highly merchant-specific,
shaped by factors such as sector, size, and transaction characteristics.
In practice, large or low-margin merchants, such as grocery stores,
often negotiate lower effective interchange rates, further complicating
efforts to generalize from posted fee schedules.

What are the main facts:

- large heterogeneity across merchants

- large heterogeneity WITHIN merchants --> sufficient statistic for
transfers

- large heterogeenity in heterogeniety

- size ameliorates effects

This paper asks two central questions: (1) how do consumers pay for
goods and services across different types of merchants, and (2) what
are the implications for merchant costs, prices, and consumer welfare?
To address these questions, we utilize novel merchant-level data from
Fiserv, which covers roughly one-fifth of U.S. card transactions,
to analyze the cost and composition of payments across a wide range
of retailers. We combine this reduced-form analysis with a quantitative
framework of retail competition to estimate the redistributive effects
of interchange fees.

The paper proceeds in three parts. First, we document striking heterogeneity
in both payment mix and interchange fees across merchants. While some
merchants predominantly receive low-cost payments, such as cash and
debit, others are heavily skewed toward high-cost premium credit cards.
We then explore the determinants of this variation and its implications
for merchant-level acceptance costs. Second, we leverage this heterogeneity,
along with a regulatory shock, to estimate the causal effect of interchange
fees on retail prices. We find that, on average, merchants pass these
fees to consumers in the form of higher prices at a rate close to
one-for-one. These first two parts are essential for understanding
redistribution and for the final section of the paper: they establish
how payment mix influences merchant costs and how those costs translate
into consumer prices, providing the empirical foundation for our quantification
of redistribution. In the third part, we embed these findings in a
sufficient-statistics framework to measure how the payment system
redistributes consumption through retail prices and card rewards.

Our results highlight the magnitude and limits of redistribution through
the payment system. We estimate that interchange fees generate approximately
\$\scalarinput{combined_dollar_headline} billion in annual transfers
from cash and debit card users to credit card users each year. These
transfers are economically significant, comparable in size (but opposite
in direction) to major government programs such as SNAP (\$120bn),
the Earned Income Tax Credit (\$57bn), and unemployment insurance
(\$50bn). Put differently, cash users pay the equivalent of a 20\%
higher sales tax than premium credit card users.

Critically, however, two mechanisms substantially attenuate this transfer.
First, consumers who use different payment methods tend to shop at
different types of merchants, reducing the scale of cross-subsidization
by approximately \scalarinput{pct_reduction}\%. Since cross-subsidization
occurs at the merchant level, there is no redistribution if cash,
debit card, and credit card users shop at different merchants. Second,
variation in negotiated interchange rates, especially the lower fees
paid by large, low-margin merchants like grocery stores, further reduces
transfers by approximately \$\scalarinput{low_income_firm_redist}
billion annually. While the direction of redistribution aligns with
theory, the degree of transfer is significantly reduced by consumer
sorting and merchant bargaining power.

To study the effect of interchange fees, we draw on two comprehensive
datasets from Fiserv, a large U.S. merchant acquirer. The first is
transaction-level settlement data from all Fiserv merchants from 2009
to 2022, covering the value, payment method, network, and interchange
fee for each card transaction. We focus primarily on a 2022 cross-section
of approximately 1 million merchants\footnote{The transaction-level data allows us to study both payment volumes
at the outlet level, as well as at the corporate level. We refer to
this latter as a merchant.}, which covers around one-fifth of all U.S. card payments. The second
dataset contains transaction data from Fiserv's Clover platform, covering
approximately 800,000 merchants between 2019 and 2022. Clover is the
largest U.S. cloud-based point-of-sale (POS) system and processed
over \$133 billion in transactions as of 2020.\footnote{Source: https://web.archive.org/web/20230717211716/https://investors.fiserv.com/static-files/820517ed-48a8-4925-aca7-181065168e25}
A unique feature of the Clover data is that it includes both cash
and card payments, allowing us to observe the full composition of
payments. Cash transactions are notoriously difficult to capture in
large datasets. We augment these sources with merchant-level characteristics
such as sector and county, and validate the data against public sources
including the IRS, the U.S. Census, and the Diary of Consumer Payment
Choice.

In the first part of the paper, we document two basic facts about
heterogeneity in the cost and composition of payments across merchants.
Our first fact is that the composition of payments---defined by the
share of payments made with a payment method---varies substantially
across merchants. Given the large fee differences between card types,
the variation in payment composition necessarily implies substantial
heterogeneity in the cost of payments across merchants. For example,
although debit card transactions account for a little more than half
of all card transactions, the share of debit transactions across merchants
has a standard deviation of approximately $24\%$. The distribution
is also bimodal: some stores have almost all debit card shoppers,
while others have few. This variation constrains the extent to which
credit card interchange fees redistribute consumption from debit card
consumers.

The composition of payments across sectors and geographies is consistent
with the fact that higher-income individuals are more likely to use
credit cards, especially premium credit cards, and less likely to
use debit cards. Travel and retail tend to feature more credit card
expenditures, whereas grocery and gas stations tend to feature more
debit cards. Within retail, luxury categories, such as jewelry stores,
are more likely to receive credit card payments, whereas shoe stores
are more likely to receive debit card payments. Consumers in higher-income
counties are also more likely to use credit cards. Ultimately, the
variation in the composition of payments, sector, and firm size explains
around \scalarinput{icg_r2}\% of the variance in average interchange
fees across merchants, which, as we discuss later, has important implications
for redistribution.

Our second fact is that card type, merchant sector, and merchant size
play fundamental roles in determining merchant-level interchange fees.
While much of the prior literature has focused on the distinction
between debit and credit, we show that different types of debit and
credit cards charge significantly different fees. For example, debit
cards issued by large banks regulated under the Durbin Amendment charge
\scalarinput{fees_nobargain_regulated_debit}\%  whereas debit cards
issued by smaller, exempt banks charge around \scalarinput{fees_nobargain_unregulated_debit}\%.
These high-fee exempt debit cards for account for 40\% of debit transactions.Within
credit cards, basic credit cards charge merchants around \scalarinput{fees_nobargain_basic_credit}\%,
while high-fee/high-rewards premium credit cards, often used by higher-income
consumers, charge around \scalarinput{fees_nobargain_premium_credit}\%.
Credit card premiumization has accelerated in recent years; premium
credit cards represented 20\% of credit card volume in 2010 and now
account for 60\% as of 2022. The fee variation across cards plays
an increasingly important role in determining merchant-level fees,
as cash use---around $12\%$ of transaction value in the Clover data---has
declined.\footnote{Payment methods in the U.S. have shifted gradually over time, with
cash use declining and debit card use steadily increasing. The use
of cash and checks fell from approximately 13\% of consumer payments
in 2019 to around 11\% in 2022. Since 2005, debit cards have gained
share relative to credit cards: while credit accounted for roughly
60\% of card payment volume in 2005, it accounted for only 50\% in
2022.}

Even holding card types fixed, interchange fees also vary across merchants.
Merchants in high-gross-margin sectors, such as travel and retail,
pay high interchange fees, whereas grocery stores and gas stations
pay around $50$ bps lower interchange fees. Large merchants in the
grocery and retail sectors, with annual card sales exceeding $\$1$
billion tend to pay lower interchange rates, which are approximately
$50$ basis points lower, compared to smaller merchants. These lower
rates are the consequence both of publicly posted volume discounts
as well as private negotiation and reflect the bargaining power of
larger merchants.

In the second part of the paper, we present novel evidence that interchange
fees have real effects on merchant sales. To do this, we exploit variation
in the effects of the $2011$ Durbin Amendment across merchants, which
has been exploited in other work such \citet{Mukharlyamov.Sarin2025}.
The Durbin Amendment capped interchange fees only on debit cards issued
by large banks, but not on debit cards issued by small banks, nor
on credit cards. It also inadvertently increased debit card interchange
fees for merchants with small average transaction sizes, while creating
significant savings for merchants with large transaction sizes. We
implement an instrumental variables design that compares sales growth
across merchants with different predicted savings from the Durbin
Amendment, using variation in local exposure to large banks to isolate
causal effects. We find evidence suggesting that Durbin Amendment
led to a decline in prices and, ultimately, find that every one percentage
point in interchange expense saved causes around a \absinput{durbin_iv_2Y_Sales_IV}\%
increase in card sales.\footnote{Using zip-code level data on retail gasoline prices, \citet{Mukharlyamov.Sarin2025}
find that they lack sufficient power to identify the effects of interchange
fees on retail prices. In our setting, we are able to use the Fiserv
data to show that merchant-level shocks to interchange fees do affect
merchant-level sales.}

In the final part of the paper, we use these facts about payment mix,
interchange fees, and merchant pricing to analyze how changes in interchange
fees affect consumer welfare. Our empirical evidence highlights the
importance of accounting for consumer sorting and interchange fees
at the merchant level. We incorporate these sources of merchant-level
heterogeneity into our analysis by developing sufficient statistics
that capture the consumer welfare effects of changes in baseline interchange
fees, sector-level pricing, and large-merchant discounts. These statistics
remain valid even when consumers' payment and shopping preferences
are arbitrarily correlated. The sufficient statistics apply to any
model in which merchants pass through interchange fees into retail
prices and operate under imperfect competition. The key insight of
our approach is that the effects of consumer sorting across merchants
can be summarized by the distribution of payment shares across merchants.

Even after accounting for the complexities of the interchange schedule
and the variation in payment composition across merchants, we find
that the payment system generates large transfers from cash and debit
card users to credit card users of around $\$30$ billion per year.
Our calculations suggest that cash users lose around \scalarinput{combined_cash_final}
basis points of purchasing power and that regulated debit card users
lose around \scalarinput{combined_regulated_debit_final} basis points.
Although much of the emphasis has been on cash users subsidizing credit
card rewards, our results suggest that debit card users also subsidize
a substantial share of those rewards. Relative to an average state
and local sales tax rate of around $6\%$, our calculation suggests
that interchange fees are analogous to raising the sales tax rate
for cash users by around 17\% and that for regulated debit card users
by 10\%. In contrast, basic and premium credit card users consume
around \scalarinput{credit_avg_cons} basis points more as a result
of high interchange fees, effectively reducing their effective sales
taxes by 10\%.

At the same time, we find that consumer sorting matters for the magnitude
of this transfer and, for specific groups, can even change the direction
of the transfer. Because cash, debit, and credit card users tend to
shop at different types of stores, consumer sorting reduces the effective
size of the transfer by approximately \scalarinput{credit_overstate}
billion per year. This effect is largely driven by the tendency of
many high-income, premium credit card users to shop at merchants that
primarily serve similar customers. If a merchant primarily accepts
premium credit cards, there is little to no cross-subsidization. As
a result, consumer sorting insulates cash and debit card users from
the full effects of the rise in premium credit card usage.

Variation in interchange fees across merchants, often driven by bargaining
power and volume discounts, also has important implications for redistribution.
Because redistribution occurs at the merchant level, it is most pronounced
at merchants where consumers exhibit substantial heterogeneity in
payment methods. Grocery stores, gas stations, and large retailers---where
all types of consumers shop---have wide variation in payment mix
and, therefore, the greatest potential for redistribution. However,
our reduced-form analysis shows that these very merchants, despite
having the most scope for redistribution, are often able to negotiate
the lowest interchange rates. Lower interchange fees inherently limit
the extent of redistribution. As a result, the fact that merchants
with the greatest scope for redistribution face the lowest fees significantly
reduces the total amount of redistribution. These negotiated discounts
lead to a progressive redistribution of approximately \$\scalarinput{low_income_firm_redist}
billion per year to cash and exempt debit card users.

We also use our framework to examine how two recent shifts in the
industry, the Durbin Amendment and the rise in card premiumization,
have impacted redistribution. Although cash users benefited from the
Durbin Amendment, our analysis shows that credit card users were the
primary beneficiaries. Our results also indicate that the Durbin Amendment
was a regressive policy, benefiting wealthier households on net at
the expense of lower-income households. Similarly, we find that the
increase in credit card premiumization has been regressive and has
altered the nature of redistribution. The rise in premiumization has
come largely at the expense of basic credit card users.

The paper proceeds as follows. Section \ref{sec:institutional-data}
describes the institutional setting and the two primary data sources,
including merchant-level transaction data from Fiserv\textquoteright s
settlement platform and the Clover point-of-sale system. Section \ref{sec:costs}
documents variation in the cost and composition of payments across
merchants, highlighting the role of card type, sector, and merchant
size. Section \ref{sec:real_effects} estimates the causal effect
of interchange fees on merchant outcomes, using variation induced
by the Durbin Amendment. Section \ref{sec:redistribution} develops
a sufficient-statistics framework to quantify redistribution across
consumers, accounting for consumer sorting and merchant bargaining.
Section \ref{sec:conclusion} concludes.

\subsection{Related Literature}

Our paper is related to several different strands of the literature.
First, our paper relates to the literature on payment choice. Much
of this literature has relied on survey evidence \citep{foster20112009,bagnall2014consumer,greene20182016,huynh2022equilibrium}.
One commonly used data source is the Survey and Diary of Consumer
Payment Choice, which typically surveys 1,000 to 5,000 consumers annually.
Our work complements the existing literature by examining high-frequency
transaction-level data. Surveys can help understand individual-level
payment preferences. However, as we demonstrate in our framework,
observing payments and the associated costs at the merchant-level,
as we do in our data, is essential for understanding cross-subsidization.

Our paper examines the impact of interchange fees on prices and the
potential implications for consumer rewards. Much of the prior evidence
on the effects of interchange fees comes from regulatory interventions.
Previous studies show that when interchange fees are capped by regulation,
merchant fees decline nearly one-for-one \citep{Gans2007,Valverde.etal2016,EuropeanCommission2020}.Caps
on interchange fees also affect consumer rewards. Following the EU\textquoteright s
2015 credit card interchange cap, UK banks scaled back credit card
rewards \citep{Gunn2016}. Similarly, after Australia capped interchange
fees on Visa and Mastercard credit cards in 2003, banks substantially
reduced their credit card reward programs\citep{Chan.etal2012}.

Our analysis of the Durbin Amendment builds on the findings of \citet{Mukharlyamov.Sarin2025},
who find that the amendment led some retailers to pass on lower prices
to consumers. While much of their work focuses on how the Durbin Amendment
affected checking accounts and banking activity more broadly, the
authors also examine its impact on gasoline prices.\footnote{\citet{manuszak2017impact} and \citet{kay2018competition} also find
that the Durbin Amendment impacted bank deposit fees and prices.} As they note, a limitation of their analysis is that they observe
interchange fees, card transactions, and gasoline prices only at the
ZIP code level and with annual frequency as of 2016. We extend their
analysis by examining monthly transactions at the merchant level during
the years immediately surrounding the Durbin Amendment. This provides
a more granular and timely assessment of its impact. We also use our
data to measure precisely how the Durbin Amendment affected merchants\textquoteright{}
costs. 

Our paper also relates to the literature examining how cash payments
potentially subsidize card payments, dating back to \citet{carlton1994antitrust}.
In related work, \citet{schuh2010gains,Felt.etal2023} develop a calibrated
model to quantify who gains and loses from credit card payments. Empirically,
while we find that card use imposes a cost on cash users, the amount
of subsidization is much smaller in our data. However, we add the
caveat that we are studying a more recent period, during which card
use has become dominant. Our findings are broadly consistent with
\citet{Wang2025}, who develops a structural model of payment competition
and find that network competition exacerbates the problem of high
merchant fees and excessive card use. In similar work, \citet{huynh2022equilibrium}
develop and estimate a structural model to understand competition
in the payments sector; the authors find that a decline in interchange
fees could be welfare enhancing.

Our paper also relates to the literature on Fintech and payments more
generally. Evidence from \citet{Higgins2024} suggests that consumers
and merchants benefit from the introduction of debit cards in Mexico.
At the same time, \citet{alvarez2020consumer} find that having a
choice over payment methods, including cash, is valuable to consumers.
Understanding the value of cash is particularly important given the
potential implications of central bank digital currency \citep{whited2022will}.


\section{Institutional Setting and Data\label{sec:institutional-data}}

In this section, we outline the institutional details and describe
our data. We begin by explaining how credit and debit card payment
systems operate, focusing on the key participants as well as the associated
costs, benefits, and transfers. We then introduce our two primary
transaction-level payments datasets from Fiserv, which allow us to
observe how the types and costs of payment methods vary across merchants.
Together, these data provide a detailed view of the U.S. payment system,
linking card types, merchant characteristics, and fees.

\subsection{Payments Processing and Interchange Fees}

A typical card payment involves the customer making a purchase, the
merchant accepting the payment, and three key financial intermediaries:
the merchant acquirer, the issuer, and the card network. The issuer,
effectively the consumer\textquoteright s bank (e.g., Chase), facilitates
the transaction from the consumer\textquoteright s perspective by
providing the card and processing the payment. Merchants, in turn,
accept card payments by contracting with a merchant acquirer, usually
a bank (e.g., Wells Fargo) or fintech firm (e.g., Fiserv, Square),
which supplies payment terminals and facilitates settlement and fund
transfers. The card network (e.g., Visa, Mastercard, American Express,
Discover) is an intermediary between the issuer and the acquirer,
routing transaction data, collecting and distributing fees, and coordinating
settlement.

To accept card payments, merchants pay transaction fees, which vary
by payment method. Figure \ref{fig:payment-flow} illustrates how
fee transfers among these parties. When a consumer makes a card payment,
the merchant pays a merchant discount fee to its acquirer. While the
acquirer retains a portion of this fee, it passes most of the fee
to the issuer as an interchange fee. Additionally, both the issuer
and the acquirer pay network fees to the card network.The interchange
fee is typically the largest component of the merchant discount fee
and funds consumer rewards, and consequently, is the focus of much
of our analysis.

- TABLE 4 DOES NOT CORRESPOND TO WHAT WE SAY IT CORRESPONDS TO--IT
HAS OUR OWN DEFINITIONS; WE NEED TO POST A PAGE OF THE ACTUAL SCHEDULE
BECAUSE THE VARIATION IS SO RICH and it shows how difficult it is
to compute fee differences

(THE TABLE DOES NOT SAY EXCERPT; IT DOES NOT SAY IN PANEL B THAT THIS
IS AN EXCERPT FOR RESTAURNATS--MAKES IT SEEM LIKE FOR ALL FIRMS IS
QUICK GLANCE) .

Table \ref{tab:visa-interchange-public} presents an excerpt from
Visa\textquoteright s 2022 public interchange schedule to highlights
how interchange fees vary with type of card, the size of the transaction,
sector of the merchant, and the size of the merchant on the example
of the restaurant industry. The schedule also shows there are many
other sources of variation, such as how the payment was collected
also influence fees paid by merchants. The table highlights that unless
a researcher knows the full set of fees paid or the exact type of
transactions, computing either the fees paid by merchants, or the
transfers associated with them will be difficult. Moreover, fees can
depart from the posted schedule due to negotiation, further complicating
the calculation. Our data captures actual fees paid by merchants,
and does not rely on posted schedules, but schedules are nevertheless
informative to understand several institutional dimensions that will
play a central role in our analysis.

\textbf{Card type classification and regulation}

As the table highlights, interchange fees vary substantially by card
type. For ease of interpretation, we classify card payments into four
categories: regulated debit, exempt debit, basic credit, and premium
credit.

We classify debit card transactions based on whether the issuing bank
is regulated or exempt under the Durbin Amendment. The 2011 Durbin
Amendment capped interchange fees for debit cards issued by banks
with over \$10 billion in assets, setting a maximum fee of 22 cents
plus 0.05\% per transaction. Banks below this threshold remained exempt
and continued to charge higher fees. For comparison, the European
Union's interchange fee regulation (IFR) caps the interchange fee
on all debit cards at $0.2\%$

We classify credit card payments into basic and premium categories.
Premium cards include those in higher interchange tiers, such as Visa
Signature, Visa Infinite, Visa Signature Preferred, Mastercard World,
and Mastercard World Elite. These cards offer enhanced rewards and
carry higher interchange fees.\footnote{Rewards can still vary within card type because the exact structure
of rewards are determined by issuers. Visa and Mastercard influence
the set of rewards primarily by changing the interchange fees issuing
banks can earn for different card types.} Basic credit cards include all other credit cards. Table \ref{tab:public-icg-sector}
shows that premium credit cards have publicly reported interchange
rates that are higher than those for basic cards. The fees on credit
cards are not regulated in the U.S. but there are several proposed
pieces of legislation at both the federal level (Credit Card Competition
Act at the federal level or Colorado\textquoteright s \textquotedblleft Swipe
Fee Fairness and Consumer Safeguards Act\textquotedblright{} at the
state level). The European Union regulation caps the interchange fee
on all credit cards at $0.3\%$ of transaction value.

\textbf{Fixed and variable fee component}

Interchange fees often contain a fixed component and a variable component
that scales with the transaction size. For example, for Visa's posted
schedule in Table \ref{tab:visa-interchange-public}, the interchange
fee on a regulated debit card at a restaurant is $22$ cents plus
$0.05\%$ of transaction value. For a $\$5$ cup of coffee, this equals
around $22$ cents, or around $4.4\%$ of transaction value. If the
same card were instead used to buy a $\$100$ dinner, the interchange
fee for the regulated debit card would only rise to $27$ cents, or
only around $0.3\%$ of transaction value. This induces variation
in interchange fees across merchants with different average transaction
sizes. As we later illustrate, because of the fixed component, regulated
debit cards can have higher costs than exempt debit card for small
transactions. We exploit this variation in Section XXX to estimate
the pass-through of fees.

\textbf{Sectors}

Interchange rates also vary across sectorof the merchant. For example,
the second row shows that exempt debit cards charge grocery stores
flat $30$ cent interchange fees. Thus, while a $\$40$ meal at a
restaurant would incur around $58$ cents of interchange, the interchange
fee on a $\$40$ basket of groceries is only around half of that.

\textbf{Merchant Size}

Interchange fees also vary with the size of the merchant. This reflects
both publicly posted volume discounts as well as private negotiations.
While the terms of negotiated fees are not posted, we observe actual
fees paid in our data. Panel \ref{tab:public-icg-firm-size} shows
the publicly posted volume discounts received by supermarkets for
premium credit cards. For supermarkets with less than $\$1$ billion
in sales, the interchange fee on a premium credit card is around $2\%$
of transaction value plus $7$ cents. In contrast, the interchange
fee for supermarkets with more than $\$20$ billion in annual sales
is around $30\%$ less, at around $1.4\%$ of transaction value plus
$2$ cents.

In principle, interchange fees can be used to compensate issuers for
fraud as well as credit risk. However, the total cost of card fraud
in the U.S. is around $10$ bps of transaction value, and merchants
bear around $80\%$ of that loss \citet{Hayashi.etal2016}. In particular,
fraud rates are much higher for online transactions, and merchants
in the U.S. bear the entirety of the cost of online card fraud. Interchange
fees are unlikely to compensate for credit risk because the set of
consumers that borrow is largely distinct from the set of consumers
who spend and thus earn issuers interchange. For example, \citet{Adams.Bord2020}
find that credit card transactors account for around $70\%$ of all
credit card transaction volume, but only $10\%$ of the interest expense.
\citet{Agarwal.etal2022} find that the probability of revolving on
a credit card decreases substantially in credit score whereas transaction
volumes increase with credit score. Thus the revenue for low FICO
borrowers is largely from interest expense and fees, and relatively
little from interchange.

\subsection{Data}

We combine data on payment volumes and interchange fees from Fiserv,
a large merchant acquirer, with data from consumer surveys to quantify
the extent the payment system redistributes consumption.

\emph{Fiserv's payment settlement records}. Fiserv is a merchant acquirer
responsible for processing transactions and calculating the interchange
fees owed to issuers. The data include both \emph{in-person} and \emph{online
transactions}. The key advantage of this data are its granularity
and breadth, which allow us to analyze how payment fees vary across
merchants. We observe payment methods at the establishment (i.e. store)
level, allowing for a detailed analysis of the distributional impacts.
We also capture fine-grained information on specific payment types,
including distinctions between regulated versus exempt debit cards,
credit card tiers, as well as cash transactions at a subset of merchants.
We derive four distinct samples from the Fiserv data for our analysis.
\begin{enumerate}
\item We construct a cross-sectional dataset of sales, the number of transactions,
and total interchange fees at the merchant-sector-location-card type
level from $2022$. This dataset covers around $1$ million merchants
and one-fourth of U.S. card volumes. We use this data to study the
composition and pricing of payments. All sales and interchange values
are measured net of returns and chargebacks. All merchants are on
card networks are assigned a Merchant Category Code (MCC). We aggregate
these MCC into six sectors -- grocery, merchandise, gas stations,
travel, restaurants, and other, and use these as our sector classifications.
The data on sectors is particularly valuable because interchange fees
vary by sector. GM NOTE: THIS LOOKS SO WEIRD-->NEED OT EXPLAIN WHY
WE USE ONLY 2022 FOR THIS RATHER THAN IT'S SOME WEIRD CHERRY PICKING.
WHY NOT THE TERMINAL YEAR OF OUR SAMPLE (2023). IF SO, THEN WE DO
NOT HAVE TO DESCRIBE IT SEPARATELY FROM THE OVERALL DATA IN POINT
4
\item We construct a three-year panel of sales, transaction counts and interchange
fees by card types at restaurants from $2010$ to $2013$. We use
this dataset to study the effect of the Durbin Amendment, a regulatory
cap on debit interchange fees. In that dataset, we aggregate the data
to the establishment-by-payment method-by-year level. The dataset
allow us to categorize merchants according to their average ticket
size (i.e. total sales divided by total transactions), which is an
important input into explaining how the effects of the Durbin Amendment
varied across merchants.
\item We construct an four-year imbalanced panel spanning of merchants that
use Clover, a Fiserv-owned point-of-sale (POS) system. This data covers
approximately 800,000 merchants and covers the period from $2019$
to $2022$. The main advantage of the Clover data is that it includes
cash and check transactions, which are typically unavailable in settlement
data. Similar to the settlement data, the Clover dataset contains
information on the value and count of transactions by payment method,
as well as data on the sector and location of the merchant. However,
it does not have data on interchange fees, and the sample is sample
skews toward smaller merchants and over-represents restaurants.
\item We construct time series of total payment volumes and interchange
fees by the type of card. These time series cover $2005-2023$ and
give context for general changes in the composition and pricing of
payments over time.
\end{enumerate}
\emph{Atlanta Federal Reserve's Diary of Consumer Payment Choice (DCPC)}
from the Atlanta FED. It records the transactions of around 5,000
respondents each year and tracks their payments over a three day window
in October. Crucially for our applications, it provides us data on
consumers' incomes and their preferred payment methods.

\emph{MRI-Simmons Ultimate Survey of Americans (USA)}. The MRI-Simmons
USA surveys around $50,000$ respondents each year on their brand
preferences, including for financial products. The data on bank choice
and credit card type allow us to group consumers by more granular
payment preferences that are relevant for the rewards they receive
and the interchange fees they impose on merchants.

\subsubsection{Data Coverage and Representativeness--NEEDS MUCH MORE}

The Fiserv data is highly representative of the broader U.S. economy
in terms of sectoral coverage, firm size (based on card sales), and
geographic distribution. Figure \ref{fig:data-representative-sector}
compares the sector composition of card transactions in the Fiserv
data to that reported in the Diary of Consumer Payment Choice. We
observe close alignment across major retail categories such as restaurants,
grocery stores, merchandise retailers, gas stations, and travel.

Figure \ref{fig:data-representative-geo} compares the state-level
distribution of card transactions in the Fiserv data with state shares
of GDP. Most points lie along the 45-degree line, indicating strong
geographic representativeness. Figure \ref{fig:data-representative-size}
compares firm size distributions in the Fiserv dataset with those
reported by the IRS Statistics of Income, again showing close correspondence
across revenue categories.


\section{Merchant Cost of Payments\label{sec:costs}}
\begin{itemize}
\item large differences in fees across merchants
\begin{itemize}
\item could be fine if all consumers shop at same place
\end{itemize}
\item sd
\end{itemize}
We begin by documenting differences in interchange fees across merchants
and within merchants. The within-merchant dispersion is especially
relevant

We then decompose the sources of this variation

merchants in payment methods and the variation in payment behavior
across merchants. We show that consumers\textquoteright{} choice of
payment method varies substantially across and within industries.
At some merchants, consumers primarily pay by credit card, while cash
and debit dominate at others. Fees also vary by the type of card,
the sector of the merchant, and the size of the merchant.

The heterogeneity in the composition and pricing of payments has two
essential implications for quantifying the amount of redistribution
in the payment system, as we illustrate in Section \ref{sec:redistribution}.
First, cross-subsidization only arises when consumers with different
payment preferences shop at the same set of merchants. Second, redistribution
is further reduced (amplified) if the firms that serve the most heterogeneous
consumer bases pay lower (higher) interchange fees.

\subsection{Large Differences across in Costs of Payments across Merchants}

We first document substantial variation in merchant transaction fees
for accepting card payments using payment settlement records in Figure
\ref{fig:interchange_dist_card} . On average, merchants pay around
\scalarinput{icg_summary_mean}\% of sales to accept cards, but this
average masks considerable variation. At the 10th percentile, interchange
fees are \scalarinput{icg_summary_p10}\% of sales, while at the 90th
percentile, they rise to \scalarinput{icg_summary_p90}\% of sales.

Part of the variation in average transaction costs is driven by the
mix of cash versus card payments across firms, as cash transactions
incur zero interchange fees, while card interchange fees can be upwards
of two percent. However, there is still substantial variation in card-related
costs across merchants. Figure \ref{fig:interchange_dist_card} displays
the average card transaction costs across merchants in our sample
with more than \sizecutoff in annual sales. The average transaction
cost is \scalarinput{icg_summary_mean}\%, with a standard deviation
of 0.5\%.

We next examine the factors driving this variation in payment costs.
The overall payment mix---including the types of debit cards (regulated
vs. exempt) and credit cards (basic vs. premium) used---has a first-order
impact on merchant costs. Additionally, even within a given payment
type, cost variation arises due to merchant-specific characteristics,
such as industry, firm size, and the average transaction size at a
merchant. We begin by documenting the variation in transaction costs
and conclude by decomposing this variation into its key components.

\subsection{Variation in Payment Costs}

\subsubsection{Card Types\label{subsec:card-types}}

We classify card payments into four categories: regulated debit, exempt
debit, basic credit, and premium credit. Distinguishing between these
card types is crucial because they impose different costs on merchants.
Figure \ref{fig:interchange_dist}, panels (a)--(d), and Table \ref{tab:icg-summary}
display average interchange rates for each of these categories. Regulated
debit cards earn around \scalarinput{fees_nobargain_regulated_debit}\%
of transaction value in interchange, while exempt debit cards earn
higher rates around \scalarinput{fees_nobargain_unregulated_debit}\%.
Basic credit cards carry higher interchange rates, around \scalarinput{fees_nobargain_basic_credit}\%,
and premium credit cards have even higher rates, averaging approximately
\scalarinput{fees_nobargain_premium_credit}\%.

We classify debit card transactions based on whether the issuing bank
is regulated or exempt under the Durbin Amendment, which capped interchange
fees for debit cards issued by large banks (regulated debit) but not
for those issued by smaller institutions (exempt debit). Specifically,
the 2011 Durbin Amendment capped interchange fees for debit cards
issued by banks with over \$10 billion in assets, setting a maximum
fee of 22 cents plus 0.05\% per transaction. Banks below this threshold
remained exempt and continued to earn higher fees. Table \ref{tab:public-icg-sector}
compares reported exempt and regulated debit interchange rates across
sectors using Visa\textquoteright s public interchange schedule. At
an average debit card transaction size of around \$40, exempt rates
are substantially higher than regulated rates in many sectors. Figures
\ref{fig:interchange_ts} and \ref{fig:interchange_debit_ts} display
average interchange fees for debit transactions overall and separately
for regulated and exempt debit cards. Figure \ref{fig:interchange_ts}
shows that average debit card interchange fees fell by more than 50\%
following the implementation of the Durbin Amendment.

Figure \ref{fig:regulated_debit_ts} shows the share of debit card
transactions that are subject to the Durbin Amendment cap over time.
By definition, the share of regulated debit transactions was zero
before the amendment, and it jumped to 60\% following its introduction
in 2011. This share has remained roughly constant since then.

We classify credit card payments into basic and premium categories.
Premium cards include those in higher interchange tiers, such as Visa
Signature, Visa Infinite, Visa Signature Preferred, Mastercard World,
and Mastercard World Elite. These cards offer enhanced rewards and
carry higher interchange fees.\footnote{Rewards can still vary within card type because the exact structure
of rewards are determined by issuers. Visa and Mastercard influence
the set of rewards primarily by changing the interchange fees issuing
banks can earn for different card types.} Basic credit cards include all other credit cards. Table \ref{tab:public-icg-sector}
shows that premium credit cards have publicly reported interchange
rates that are higher than those for basic cards. Figures \ref{fig:interchange_ts}
and \ref{fig:interchange_credit_ts} display average interchange fees
for credit card transactions overall and separately for basic and
premium credit cards. Figure \ref{fig:interchange_ts} shows that
average credit card fees have increased over time. However, this increase
is primarily driven by the growing prevalence of premium credit cards,
as the average interchange fees for basic and premium cards have remained
relatively stable over this period (Figure \ref{fig:interchange_ts}).
Figure \ref{fig:premium_credit_ts} illustrates the rising dominance
of premium cards, which grew from about 20\% of credit card volume
in 2006 to 60\% in 2022.\footnote{Historically, card networks such as Visa and Mastercard have required
that merchants accepting their cards must accept all cards issued
under their respective brands. In other words, merchants cannot selectively
accept certain Visa cards---they must accept all Visa cards or none
at all. This policy has been the subject of ongoing \textquotedblleft honor-all-cards\textquotedblright{}
litigation.}

We examine the extent to which differences in payment mix can explain
variations in transaction costs. Specifically, we predict each merchant\textquoteright s
average transaction costs based on their payment mix, including variations
within card types (e.g., regulated vs. exempt debit and basic vs.
premium credit). Variation in payment mix accounts for approximately
\scalarinput{icg_r2_shares}\% of the overall variation in transaction
costs. The dashed red line in Figure \label{fig:interchange_decomp}
shows the density of predicted transaction costs based on each merchant\textquoteright s
observed payment mix. The remaining variation in transaction costs
is attributable to within--card-type differences in interchange fees,
which reflect variation across merchants due to factors such as industry,
firm size, and transaction size.


\subsubsection{Merchant Sectors}

Interchange rates also vary across sectors, with high-margin sectors
typically paying higher fees. Figure \ref{fig:price-disc-merch} plots
average interchange rates in the Fiserv data across five major spending
categories: travel, restaurants, merchandise, grocery stores, and
gas stations. These five categories account for approximately \scalarinput{share_sector}\%
of card sales. We focus on rates for credit cards and exempt debit
cards, as the Durbin Amendment cap limits variation in regulated debit
card fees.

Fees are highest in the travel and retail sectors, which generally
have higher margins and more affluent customers compared to other
sectors. In contrast, grocery stores and gas stations, which operate
on thinner margins and serve a broader swath of the income distribution,
pay interchange fees that are one-third or \scalarinput{grocery_discount}
basis points lower. One natural explanation for why high margin merchants
are more willing to accept cards is that interchange fees are charged
as a percentage of sales. Even if accepting credit cards has the same
impact on sales across sectors, card acceptance would increase profits
at high-margin businesses by a greater amount . Historically, Visa
and Mastercard had lower acceptance at supermarkets until they started
to give supermarkets discounts on interchange rates in the 1990's
\citep{Stearns2011}.

Accounting for the merchant sector, in addition to the payment mix,
explains roughly \scalarinput{icg_r2_sector}\% of the overall variation
in transaction costs. The dashed green line in Figure \ref{fig:interchange_decomp}
shows the distribution of predicted transaction costs based on each
merchant\textquoteright s observed payment mix and sector.

\subsubsection{Merchant Size and Other Characteristics}

Interchange fees are generally lower for larger firms. Figure \ref{fig:price-disc-size}
plots average interchange rates in the merchandise and grocery store
sectors for large firms (with more than \$1 billion in sales) and
smaller firms (with less than \$1 billion in sales). Larger firms
pay approximately \scalarinput{bargain_discount} basis points less
in fees on their credit card and exempt debit card transactions. This
pattern aligns with public data from Visa's interchange schedule,
shown in Table \ref{tab:public-icg-firm-size}.

One explanation for this size-based discount is that large merchants
can influence consumer payment behavior by choosing not to accept
specific payment methods. In response, networks offer lower interchange
rates to large merchants to keep them on the platform. Historically,
large department stores were slow to accept Visa and MC credit cards
because they had their own alternatives in the form of store credit.
Only after Visa made concessions did JC Penney's become the first
large department store to accept credit cards, and Sears held out
until the mid 1990's \citep{Nocera1995}.

Another important factor influencing the average interchange fee is
the transaction size. Visa's public interchange schedules, as shown
in Table \ref{tab:visa-interchange-public}, indicate that interchange
fee schedules include a fixed component that does not vary with the
transaction size, as well as a variable component that scales linearly
with the transaction size.

Accounting for firm characteristics such as total sales and the average
transaction size, in addition to payment mix and sector, explains
approximately \scalarinput{icg_r2_full}\% of the overall variation
in transaction costs. The dashed orange line in Figure \ref{fig:interchange_decomp}
shows the distribution of predicted transaction costs based on each
merchant\textquoteright s observed payment mix, sector, and firm characteristics.

\subsection{Variation in Payment Mix\label{subsec:payment-mix}}

One notable feature of the data is that the composition of payments
varies significantly across merchants. We examine how the use of cash,
debit cards, and credit cards differs across the merchants in our
sample and then investigate how the payment mix varies across sectors,
regions, and other firm-specific factors.

We begin by documenting heterogeneity in payment types across merchants
in Figure \ref{fig:pmt_merchant}, panels (a)-(d).

\subsubsection{Large Variation in Cash Share across Merchants}

Figure \ref{fig:cash_share} displays the distribution of the share
of cash payments across merchants, focusing on those merchants where
cash accounts for at least 2\% of sales. Observations are at the merchant-by-month
level using our Clover data, which includes information on cash payments.
We measure the cash share based on the dollar-weighted share of cash
and check payments relative to the total transaction value. We group
cash and check payments into a single category, labeled as 'cash,'
because they are distinct from card payments in that they do not incur
interchange fees for the merchant and do not generate explicit rewards
for the user. On average, cash accounts for 11\% of transactions at
a given merchant (Figure \ref{fig:cash_share_ts}). However, this
masks substantial heterogeneity. For approximately two-thirds of merchants,
cash accounts for less than 2\% of transactions. In contrast, for
the one-third of merchants where cash represents at least 2\% of sales,
it accounts for more than 30\% of their transactions---and for merchants
in the 90th percentile, over 80\%.

\subsubsection{Large Variation in Types of Cards across Merchants}

This heterogeneity is critical when considering redistribution. For
the two-thirds of firms where cash use is almost nonexistent, there
is little scope for redistribution between cash and card users. Similarly,
at firms where cash is the predominant payment method, redistribution
is limited. Redistribution inherently requires a mixed payment environment.

\paragraph{Credit vs. Debit}

The remaining 89\% of sales occur via debit and credit card payments.
Here, too, the payment mix varies across merchants. Figure \ref{fig:credit_share}
displays the share of credit-card payments relative to total card
payments (i.e., debit plus credit) in card settlement data. Observations
are at the merchant-by-month level. The figure illustrates that, on
average, credit card sales account for 53\% of card transactions at
the typical merchant in our sample. Moreover, the share of credit
card transactions has remained relatively constant over time (Figure
\ref{fig:credit_share_ts}).

The distribution of credit card sales is bimodal, with peaks at approximately
25\% and 70\%. This bimodality suggests that although credit cards
account for 53\% of card transactions on average, merchants tend to
fall into two distinct groups: those where credit accounts for around
25\% of credit card transactions and those where it accounts for about
75\%. This pattern has implications for cross-subsidization. Debit
cards carry lower interchange fees and rewards relative to credit
cards. Just as cash transactions can subsidize credit card rewards,
debit card transactions also potentially play a subsidizing role.
The bimodal nature of the distribution suggests that variation in
card payment mix is significant at the merchant level---again limiting
the potential for cross-subsidization.

\subsubsection{Variation Across and Within Sectors}

Part of this heterogeneity in the merchant payment mix stems from
sector-specific characteristics. Figure \ref{fig:pmt_sectors} shows
the payment shares, weighted by transaction value, across sectors.
Although cash accounts for only 11\% of transactions at the average
merchant, it plays a more significant role in sectors such as grocery,
restaurants, and gas stations. For instance, cash accounts for 23\%
of sales in the grocery sector, slightly less than the share for credit
cards. By contrast, in the travel sector, credit cards account for
70\% of sales, with debit and cash representing 27\% and 3\%, respectively.
These summary statistics suggest that, in the travel sector, there
is limited scope for redistribution between credit and cash users,
and any residual redistribution is likely minimal and occurs primarily
between credit and debit users.

In addition to variation across sectors, we also find substantial
dispersion in the composition of payments within sectors. Figure \ref{fig:sector_densities}
plots kernel densities of the share of card transactions on credit
cards across merchants in each of the major sectors. Although grocery
stores tend to have a relatively low share of credit card transactions,
and travel merchants tend to have a higher share, there remains substantial
dispersion within these sectors. This dispersion is important, as
it implies that even within a given sector, consumers with different
payment preferences shop at different merchants.

\subsubsection{Variation Across Regions and Income Levels}

We also find substantial regional variation in the payment mix. Figure
\ref{fig:cash_map} displays the county-level share of cash payments.
Light green and yellow regions correspond to higher cash usage, while
dark blue areas reflect lower usage. In some counties (highlighted
in yellow), cash accounts for more than 23\% of transactions. Cash
usage is higher in the Midwest and the South, while card usage is
more prevalent along the East and West Coasts, as well as in metropolitan
areas.

Figure \ref{fig:credit_map} displays the county-level share of credit
card transactions relative to all card payments. Areas shaded in yellow
and light green correspond to higher credit card usage, whereas areas
shaded in dark green and blue reflect higher debit card usage. The
results indicate that debit card usage is especially prevalent in
the Deep South relative to other U.S. regions. Conversely, credit
card usage is higher in the coastal areas.

The geographic correlation between income and credit card usage is
also borne out in individual demographic data. Figure \ref{fig:payment_income}
shows how consumers preferred payment methods change with income.
At low levels of income, consumers prefer using debit cards and cash.
At higher levels of income, consumers use more credit cards, and in
particular premium credit cards. Thus transfers from cash and debit
card consumers are also transfers from relatively lower income consumers
to higher income consumers.

\subsubsection{Decomposing the Variation in Payment Methods}

To better understand what drives the variation in payment methods,
we estimate regressions of the form:

\begin{equation}
Cash\,Share_{jt}=X_{jt}'\beta+\mu_{l(j)}+\mu_{c(i)}+\mu_{j(i)}+\epsilon_{jt}.\label{eq:payment_type}
\end{equation}
Observations are at the merchant $(j)$-by-month $(t)$ level. The
dependent variable $Cash\,Share_{jt}$ measures the dollar-weighted
share of cash transactions relative to other payment methods. We control
for merchant characteristics in the vector $X_{jt},$ including average
transaction size and total transaction volume. We also include fixed
effects for industry $(l)$, time ($t$), and county $(c)$ fixed
effects. In addition to studying cash share, we also analyze how the
share of credit versus debit card transactions varies across merchants.

We report the corresponding estimates in Table \ref{tab:payment_type}.
Columns (1)--(3) examine cash and check as a share of all payment
volumes in the Clover data, while columns (4)--(6) analyze the share
of credit card usage within card transactions. In all columns, the
dependent variable is scaled by 100, so coefficients are interpreted
in percentage points.

The results in Table \ref{tab:payment_type} show that consumers are
more likely to use cash at merchants with higher sales volumes but
lower average transaction values. In column (1), a 100 log-point increase
in total sales is associated with a 1.37 percentage point (12\%) increase
in the share of cash transactions. In contrast, a 100 log-point increase
in average transaction value is associated with a 5.55 percentage
point (50\%) decline in cash usage. Unsurprisingly, cash usage is
negatively correlated with e-commerce: a 10 percentage point increase
in a firm\textquoteright s e-commerce share is associated with a 0.38
percentage point decline in cash usage.

Consistent with Figure \ref{fig:cash_map}, local demographics also
help explain cash usage. Consumers are more likely to pay in cash
in less wealthy, less educated, and less densely populated regions.
Column (2) shows that a one standard deviation increase in median
income is associated with a 0.14 percentage point decline in cash
usage. Similarly, a one standard deviation increase in the share of
college-educated households corresponds to a 1.07 percentage point
decline in cash usage.

In contrast, consumers are more likely to use credit cards (relative
to debit cards) at merchants with lower sales volumes and higher average
transaction values. Column (4) indicates that a 100 log-point increase
in transaction value is associated with a 9.11 percentage point (17\%)
increase in credit card usage. As with cash, demographics also influence
card mix: credit card usage is more prevalent in wealthier, older,
more educated, and more urban areas.

Overall, these results suggest that the payment mix varies substantially
across merchants, depending on characteristics such as sector and
region. As shown below, this heterogeneity has important implications
for both the cost and redistributive aspects of payments.

\section{Pass-through Of Payments Costs\label{sec:real_effects}}

In this section, we study the extent to which costs of payments are
passed though to retail prices and thus to firm sales. Our prior results
show large differences in the payment mix and associated interchange
fees, both within and across payment methods. These results suggest
that payments result in cross-subsidization between parties in the
payment system. If payment costs are fully passed though by merchants,
then consumers of payments with low interchange fees subsidize consumers
with premium payment methods. Alternatively, if merchants do not pass
through interchange fees into retail prices, then cross-subsidization
between cash and card users does not occur; instead, merchants effectively
bear the burden of subsidizing card rewards themselves. Using our
detailed merchant-level data, we exploit heterogeneity in the impact
of the Durbin Amendment across merchants to show that interchange
fees affect the allocation of consumption and retail prices.

\subsection{Differential Impact of the Durbin Amendment }

The 2011 Durbin Amendment, passed as part of the Dodd--Frank Act,
imposed a binding cap on debit card interchange fees for banks with
over \$10 billion in assets. It limited fees to a maximum of 22 cents
plus 0.05\% of the transaction amount. The rule applied only debit
cards issued by large banks (regulated debit) and left fees for debit
cards issued by small banks and credit unions unchanged (exempt debit).
As illustrated in Figure \ref{fig:interchange_ts}, the Durbin Amendment
had a substantial impact on debit interchange fees while having no
effect on credit interchange. 

The Durbin Amendment created variation in interchange fees across
merchants, driven by three forces which aid our identification. First,
only merchants with substantial sales via debit cards saw changes
in costs. Second, only merchants in areas with a high concentration
of banks subject to Durbin (large banks) experienced changes to their
fee structure. Third, because the cap includes a fixed 22-cent component,
the percentage reduction in fees increased with average transaction
size. In fact, merchants with the smallest transaction sizes experienced
an \emph{increase} in fees when facing regulated debit cards.\footnote{An inspection of Visa's public interchange fee schedules before and
after the policy suggest that most regulated debit fees bunched at
the cap. In principle, the Durbin Amendment could have forced the
networks to cut debit fees for the largest merchants even more to
incentivize said merchants to stay on the Visa and Mastercard networks.
However, Figure \ref{fig:price-disc-size} shows regulated debit fees
are similar at both large and small firms, suggesting this additional
negotiation channel was not a major force.} We use our detailed merchant-level data to measure these forces precisely.
We focus on restaurants because the variation induce by the Durbin
amendment is largest.\footnote{We exclude supermarkets and gas stations because Table \ref{tab:visa-interchange-public}
shows that those sectors saw less dramatic increases in fixed fees.
For example, gas stations went from being charged $\min\{0.75\%+0.17,0.95\}$
to $0.05\%+0.22$. Supermarkets went from being charged a flat fee
of $\$0.25$ to $0.05\%+0.22$. Over transaction sizes from $5$ to
$100$ dollars, the average change in fee for restaurants covers $[-100,120]$
bps. The same numbers for gas stations would imply half as much variation,
spanning changes in $\left[-30,65\right]$ bps. We do not focus on
large merchandise retailers because past work shows that they charge
uniform prices across establishments that face different payment costs
\citep{DellaVigna.Gentzkow2019}.}

Figure \ref{fig:durbin-het-ticket} illustrates that the Durbin Amendment
benefited restaurants with large average transaction sizes and increased
interchange fees for restaurants with small transaction sizes. This
occurred because the Durbin Amendment changed the structure of debit
interchange from a large linear fee with a low fixed fee ($1.19\%$
of transaction value + $10$ cents--WHERE DO THESE NUMBERS COME FROM)
to a low linear fee with a high fixed fee ($0.05\%$ of transaction
value + $22$ cents). As average transaction size increases, the average
interchange rate falls because the fixed component of the fee accounts
for a smaller share of total costs. Below \$XXX the large fixed component
dominates, and interchange rates actually increase for regulated debit.
In the data, the median merchant has an average ticket size of about
\$15 and therefore benefited very modestly from the Durbin Amendment.
These modest gains for the median restaurant are small relative to
the large variation in benefits across establishments.

\textbf{The vertical line marks the theoretical transaction size at
which a merchant would have been indifferent between the two fee structures;
it also roughly coincides with the point at which the two curves cross,
underscoring the value of the settlement data.} I DO NOT UNDERSTAND
THIS STATEMENT AND DO NOT SEE THE BENEFIT; IT IS LIKELY WE SHOULD
DELETE VERTICAL LINE. \footnote{We compute this by solving $0.0119x+0.10=0.0005x+0.22$. The two curves
do not cross exactly there because we approximate the average interchange
fee with the interchange fee at the average transaction size. We do
this because we do not have data on the distribution of transaction
sizes.} 

Figure \ref{fig:durbin-het-regulated} shows that the effect of the
Durbin amendment induces variation across space: after the Durbin
Amendment, average interchange rates declined sharply in areas with
a higher share of regulated banks.  In the next section, we exploit
these two dimensions of heterogeneity to examine how interchange fees
influenced merchants\textquoteright{} sales and prices.

\subsection{Pass-though estimation}

We estimate the pass-though of interchange fee to prices and the resulting
change in sales exploiting the variation induced by the Durbin Amendment.
Because interchange fees depend on merchant characteristics, which
are themselves related to sales we use a simulated IV approach: we
construct fee changes each merchant would face, based only on their
pre-policy characteristics and the regulated post-Durbin schedule.
Formally, we estimate the impact of interchange on sales and prices
using the following following specification using the restaurant industry,
which faced the largest variation in fees following the regulatory
change: 
\begin{equation}
\Delta y_{j}=\beta\Delta\tau_{j}+\epsilon_{j}\label{eq:durbin}
\end{equation}
The dependent variable, $\Delta y_{j}$, is the change in sales or
prices at merchant $j$ in the year following the implementation of
the Durbin Amendment (i.e., October 2011 to October 2012). The independent
variable of interest, $\Delta\tau_{j}$, measures the average change
in credit and debit interchange fees at the merchant over the same
period. Observations are at the merchant level.

As discussed in Section \ref{sec:costs}, interchange fees depend
on merchant characteristics, including sales volume and average transaction
size. For example, if the firms with the fastest sales growth are
also the ones that expand the debit card share of sales or increase
average transaction sizes, that would mechanically generate a negative
correlation between lower interchange fees and stronger sales growth.
The fact that most of the variation in interchange fees stems from
the Durbin Amendment, an arguably exogenous policy shock, helps alleviate
these concerns.

To address the potential endogeneity issue, we construct an instrument
based on the expected decline in interchange fees resulting from the
Durbin Amendment, derived from each merchant\textquoteright s pre-policy
characteristics. We calculate the expected change in interchange fees
as:
\[
\widetilde{\Delta\tau_{j}}=r_{z(j)}\times d_{j}\times\left(\tau^{\text{Reg Debit,Post}}(t_{j})-\tau^{\text{Debit,Pre}}(t_{j})\right)
\]
Here, $r_{z(j)}$ is the share of debit cards in merchant $j$'s zip
code that are subject to Durbin regulation; $d_{j}$ is the share
of the merchant\textquoteright s card sales made with debit cards
in the pre-Durbin period; and $t_{j}$ is the merchant\textquoteright s
average transaction size in the pre-Durbin period. The term $\tau^{\text{RegDebit,Post}}(t_{j})$
is the expected interchange rate under the post-Durbin regulated debit
fee schedule, given transaction size $t_{j}$, and $\tau^{\text{Debit,Pre}}(t_{j})$
is the expected rate under the pre-Durbin fee schedule for the same
$t_{j}$. This instrument captures the theoretical savings for a merchant,
assuming no change in transaction size or payment composition after
the policy. It exploits heterogeneity in both the expected share of
regulated debit transactions and the change in regulated-debit interchange
fees.

Our estimates suggest that the increases in interchange fees result
in a substantial decline in sales because restaurants pass almost
fully pass-though fees to prices. The IV estimates (Table \ref{tab:durbin-iv}
column 2) suggests that a one-percentage-point increase in average
interchange fees causes a \absinput{durbin_iv_1Y_Sales_IV}\% decrease
in sales. The OLS are substantially larger, consistent with the idea
that higher interchange rates are mechanically associated with lower
sales. As we Discuss in Section XXX, for example, fees decline with
firm sales and transaction size (ANY BETTER EXAMPLES). This decline
in sales could reflect either a price increase and/or a decline in
product quality (e.g., a rise in the quality-adjusted price). We proxy
for average restaurant price using the average transaction sizes.
The idea is that consumers respond to higher restaurant prices by
going to a restaurant less often while ordering the same dishes for
any given trip. IV estimates finds that a one percentage point increase
in interchange causes a roughly one-percent increase in price over
one and two year horizons (Table \ref{tab:durbin-iv} column 7 and
8). Note that OLS estimates are negative--this reflects the mechanical
negative correlation between transaction size and interchange rates
that arises when interchange fees have fixed components that do not
scale with the transaction size.

A one-for-one passthrough of costs to prices may be a bit surprising,
because the restaurants are differentiated and may have some local
market power. Even in models with imperfect competition, industry-wide
shocks exhibit high (near one-for-one) pass-through. Intuitively,
since all firms\textquoteright{} costs move together, there is little
\emph{relative }price discipline allowing firms to adjust prices almost
fully with costs. This intuition differs significant from the standard
limited pass-though of cost shocks which affect individual firms.
In that situation rivals\textquoteright{} unchanged prices provide
consumers with substitutes, so the shocked firm must absorb much of
the cost to avoid losing market share. The Durbin Amendment is closer
to an industry-wide cost shock resulting in large pass-though. Our
findings are also consistent with \citet{Amiti.etal2019} that small
firms tend to pass through costs one-for-one to prices. Assuming full
pass-through of interchange fees to prices, these estimates imply
a demand elasticity of around $8$ (IN WHAT MODEL). This elasticity
is high relative to the elasticity of around $5$ found in \citet{Sullivan2024}'s
study of restaurants using detailed transaction data. However, the
elasticities are comparable \citet{Edmond.etal2022} macro calibration
of markups.

One alternative explanation for the change in card sales we observe
is that merchants deviate from uniform pricing of payments and impose
cash-discounts or surcharges. We re-estimate our Durbin-Amendment
effects using only the ten states that prohibit any form of card surcharging
(Table \ref{tab:durbin-surcharging}).\footnote{These states are California, Colorado, Connecticut,Florida, Kansas,
Massachusetts, Maine, New York, Oklahoma, and Texas.} If changes in surcharges were driving our results we would expect
smaller or no results in the states in which surcharges are illegal.
Instead, the instrumented estimates in this no-surcharge subsample
remain virtually unchanged---and statistically indistinguishable---from
our full-sample results. This persistence demonstrates that the increase
in card sales following lower interchange fees cannot be attributed
to merchants manipulating payment choices, but rather reflects genuine
adjustments to posted prices that apply to all customers.

\section{Sufficient Statistic for Redistribution in Payments\label{sec:redistribution}}

We develop a sufficient-statistics framework to quantify the redistributive
effects of interchange fees across consumers. Our approach takes as
inputs the facts we document above: the variation in the composition
and cost of payments across merchants as well as pass-though of interchage
fees to prices. We compute the welfare effects of changes in interchange
fees when merchants can adjust retail prices in response to interchange
fees. Critically, consumers can re-optimize consumption decisions
in response to retail prices and rewards. The framework captures the
key insight that redistribution occurs only when consumers using different
payment methods shop at the same set of merchants, and that the amount
of redistribution depends on the differences in costs of payments
wihtin an individual merchants. We use this framework to quantify
the extent of redistribution generated by the payment system in the
U.S.. We evaluate the degree to which these transfers are regressive
by evaluating the impact across the income distribution. We study
how sorting across merchants mitigates this redistribution, and how
the presence of negotiated fees affects consumer welfare. We then
use the model to quantify the changes in consumer welfare if the U.S.
were to switch to EU regulation of interchange fees. 

We model how merchants, card issuers, and consumers interact in a
payment system in which interchange fees are used to fund rewards,
and joinlty they determine equilibrium retail prices and consumer
choices. The central departure from a standard economy is that card
transactions impose fees on merchants, which influence how they set
prices, while issuers recycle these fees back to consumers in the
form of rewards. 

Our framework allows for rich heterogeneity in how consumption preferences
differ across consumers with different payment preferences. For instance,
premium card users may prefer luxury retailers, while debit users
may prefer mass-market stores. The framework also allows for premium
credit card users to have different demand elasticities and substitution
patterns compared to debit card consumers.

Allow for merchants to differ in fees that are negotiated and ...

\subsection{Theoretical Framework}

The economy is a payment system where merchants consumers, and issuers
interact through prices, interchange fees, and rewards. There are
$j=1,\dots,J$ merchants in the economy, split into a set of sectors
$S$ charging price $p_{j}$. Consumers purchase goods using payments
indexed by $k=0,1,\dots,K$ provided by issuers. Payment methods earn
rewards $f_{k}$ which are funded by fees imposed on merchants $\tau_{jk}$. 

\subsubsection{Consumers}

Consumers are indexed by the payment method they use. Each consumer
with card $k$ chooses quantities $q_{jk}$ across merchants to maximize
a differentiable, strictly monotonic, and strictly quasi-concave utility
function $U$ subject to a budget constraint. They pay consumers lump-sum
rewards $f_{k}$ The budget constraint reflects available income plus
rewards.

\begin{align*}
q_{k}^{*} & =\argmax_{q_{k}}U_{k}\left(q_{1k},\dots,q_{Nk}\right)\\
\text{s.t.} & \sum_{j}q_{jk}p_{j}\leq1+f_{k}
\end{align*}
\footnote{Our assumptions rule out utility functions that treat goods as perfect
substitutes, but are compatible with nested CES and random-coefficients
logit.}

Our framework allows for rich heterogeneity in how consumption preferences
differ across consumers with different payment preferences. For instance,
premium card users may prefer luxury retailers, while debit users
may prefer mass-market stores. The framework also allows for premium
credit card users to have different demand elasticities and substitution
patterns compared to debit card consumers.

Players:

- merchants maximize profits setting prices

- a subset of merchants take fees as given 

- a subset of merchants negotates fees

- consumers maximize utility given prices and rewards

- Issuers charge interchange fees and pay rewards

\subsubsection{Modeling Interchange Fees}

We model interchange fees as linear fees that depend on the card type
and merchant. Merchant $j$ is in sector $s_{j}$, and let $b_{j}$
be an indicator equal to $1$ if merchant $j$ has a negotiated interchange
rate. Motivated by our analysis from Section \ref{sec:costs} and
public interchange schedules, the interchange fee $\tau_{jk}$ charged
to merchant $j$ for transactions made with card $k=0,1,\dots,K$
is written as:
\begin{equation}
\tau_{jk}=\overline{\tau}_{k}+\tau_{s_{j}k}+b_{j}\delta_{k}\label{eq:interchange-fee-model}
\end{equation}
where $\overline{\tau}_{k}$ is a baseline fee for card $k$, $\tau_{s_{j}k}$
is a sector-specific adjustment, and $\delta_{k}<0$ is a discount
negotiated by firms with bargaining power. We let $k=0$ denote cash.
Without any loss in generality, we normalize the sector adjustments
$\tau_{sk}$ so that the total dollar weighted value of the sector
level adjustments is zero for card $k$.

\subsubsection{Retail Pricing}

Merchant $j$ sets a uniform price $p_{j}$ for all consumers. 

GM COMMENT: I THINK IT WOULD BE BETTER TO BE GENERIC HERE. WHAT I
DO NOT LIKE HERE IS THAT REWARDS AFFECT DEMAND, BUT THAT HAS NO IMPACT
ON PRICES. 

We denote $q_{jk}^{*}$ as the equilibrium amount that consumers with
card $k$ purchase from merchant $j$. Total equilibrium sales by
merchant $j$ are $Q_{j}^{*}=\sum_{k}q_{jk}^{*}$. We assume merchant
prices are log-linear in interchange fees, given by:

\begin{equation}
\log p_{j}=\log\overline{p}_{j}+\sum_{k}\mu_{jk}\tau_{jk},\label{eq:retail-price-pass}
\end{equation}
where $\mu_{jk}=q_{jk}^{*}/Q_{j}^{*}$ represents the share of merchant
$j$'s sales from consumers using card $k$, and $\overline{p}$ is
a baseline price that reflects the merchant's marginal costs. This
linear relationship is a first-order approximation if merchants are
monopolistically competitive and consumers have CES preferences over
merchants. It also holds under logit demand when the number of firms
is large and when the markup is small relative to the marginal cost.
Importantly, this pricing assumption is also consistent with our findings
from Section \ref{sec:real_effects}, where we find that a one percent
point increase in the interchange rate causes a one percentage point
increase in prices.


\subsubsection{Issuers: Setting }

Card issuers earn zero profits. They pay consumers lump-sum rewards
$f_{k}$ that equal average interchange fees collected from merchants,
less average costs. We interpret rewards broadly as subsidies consumers
receive for card use, ranging from cash rewards on credit cards to
free checking accounts for debit cards\citep{kay2018competition,Mukharlyamov.Sarin2025}. 

\[
f_{k}=\frac{\sum_{j}p_{j}q_{jk}^{*}\left(\tau_{jk}-c_{k}\right)}{\sum_{m}p_{m}q_{mk}^{*}}
\]

WHAT IS THE BENEFIT OF EXPRESSION BELOW; MU\_JK IS UNDEFINEDDenote
merchant $j$'s equilibrium revenue by $R_{j}^{*}=p_{j}Q_{j}^{*}$,
normalized so that $\sum_{j}R_{j}^{*}=1$. Since revenues add up
to $1$, $R_{j}$ also defines a measure on the set of firms. This
allows us to express revenue-weighted sums as expectations of the
summand. Thus the per-capita reward for users of card $k$, denoted
$f_{k},$ equals average interchange fees collected on card $k$,
net of some cost:

\[
f_{k}=\mathbb{E}_{R}\left[\mu_{jk}\right]^{-1}\mathbb{E}_{R}\left[\mu_{jk}\left(\tau_{jk}-c_{k}\right)\right]
\]


\subsection{Sufficient Statistic}

We derive the sufficient statistic, which captures the money-metric
welfare effects of changes in interchange rates. The intuition behind
our result builds on Roy's identity: even though interchange fees
can affect how consumers allocate their consumption across merchants,
this margin of adjustment does not have a direct effects on consumer
welfare. Denote $V_{k}\left(p,f_{k}\right)$ as the indirect utility
function from this consumer problem, and let $\lambda_{k}$ be the
Lagrange multiplier on the budget constraint. By Roy's identity,\footnote{The regularity assumptions (WHICH) ensure a unique interior optimal
choice of consumption for every price vector $p$, and thus we can
use Roy's identity . } we have that
\[
\frac{\partial V_{k}}{\partial p_{j}}=-q_{jk}^{*}\times\lambda_{k}
\]
where $q_{jk}^{*}$ reflects consumption at the optimum. 

We are interested in how changes in interchange fees affect consumer
welfare. Let $x_{l}$ be any parameter that affects the fees on payment
method $l$ (NEED EXAMPLE). We would like to compute the money-metric
effect of changing $x_{l}$ on the welfare of consumers who use payment
method $k$. By applying Roy's identity and dropping all higher order
terms in $\tau$, we obtain

\begin{align}
\frac{1}{\lambda_{k}}\frac{dV_{k}}{dx_{l}} & =\sum_{j}-q_{jk}^{*}\frac{\partial p_{j}}{\partial x_{l}}+\frac{\partial f_{k}}{\partial x_{l}}\nonumber \\
 & =\sum_{j}-q_{jk}^{*}p_{j}\left(\mu_{jl}\frac{\partial\tau_{jl}}{\partial x_{l}}+\frac{\partial\mu_{jl}}{\partial x_{l}}\tau_{jl}\right)+\frac{\partial f_{k}}{\partial x_{l}}\nonumber \\
 & \approx\mathbb{E}_{R}\left[\mu_{jk}\right]^{-1}\left(-\mathbb{E}_{R}\left[\mu_{jk}\mu_{jl}\frac{\partial\tau_{jl}}{\partial x_{l}}\right]+I\left\{ k=l\right\} \times\mathbb{E}_{R}\left[\mu_{jk}\frac{\partial\tau_{jk}}{\partial x_{k}}\right]\right)\label{eq:sufficient-stat-formula}
\end{align}

The above expression captures the money-metric effect on a representative
consumer of payment method $k$. To convert to dollar values, we multiply
equation \ref{eq:sufficient-stat-formula} by $\mathbb{E}_{R}\left[\mu_{jk}\right]$,
the total amount of spending by payment method $k$ consumers and
the dollar amount of spending across all payment methods.

\subsubsection{Applying the Sufficient Statistic to Changes in the Interchange Fee
Schedule}

Our sufficient statistic provides intuition on the effects of changing
the overall level of fees and reducing fee heterogeneity.

\paragraph{Baseline Fees}

We model changing the baseline fee by perturbing $\overline{\tau}_{l}$
(WHAT IS THIS). Because $\frac{\partial\tau_{jl}}{\partial\overline{\tau}_{l}}=1$,
the re-distributive effect of reducing the baseline interchange level
by $\Delta$ in dollar terms for all consumers who use payment method
$k$ is:
\begin{equation}
\Delta\times\left(\mathbb{E}_{R}\left[\mu_{jk}\mu_{jl}\right]-I\left\{ k=l\right\} \mathbb{E}_{R}\left[\mu_{jk}\right]\right)\label{eq:baseline-effect}
\end{equation}
This first term of the formula says that a decrease in the overall
fees on card $l$ benefits consumers of card $k$ by lowering retail
prices by the extent to which card $k$ and card $l$ consumers overlap.
The second term then says that card $l$ consumers are hurt from lowering
fees on card $l$ because they receive fewer rewards.

The appearance of second moments makes the framework rich enough to
nest leading cases in the literature. When consumers with different
cards shop at disjoint sets of merchants, as in \citet{Gans2018},
$\mu_{jk}\mu_{jl}=0$ for all $j$ whenever $k\neq l$ and there is
no redistribution. Conversely, when merchants are homogeneous, so
that $\mu_{jk}=\overline{\mu}_{k}$ for all merchants $j$, the redistribution
depends on the aggregate share of each card type, consistent with
the intuition in \citet{Felt.etal2023}. To the extent that there
is a large mass of cash-only merchants that are both unobserved in
our dataset and not well captured in the GDP statistics, the second
moments are zero and thus ignoring them has no effect on our estimates
of redistribution. Thus, the sufficient-statistics framework captures
a wide range of economic environments---from perfect merchant homogeneity
to full segmentation---allowing us to quantify redistribution in
the presence of observed heterogeneity in payment behavior.

\paragraph{Sector Discounts}

We model eliminating the sector discounts by perturbing $\tau_{rl}$.
The dollar effect of removing all of the sector level discounts on
card $l$ is:
\begin{align}
\sum_{r}\tau_{rl}\times\frac{\partial V_{k}}{\partial\tau_{rl}} & =\sum_{r}\mathbb{E}_{R}\left[\mu_{jk}\mu_{jl}I\left\{ s_{j}=r\right\} \right]\tau_{rl}-I\left\{ k=l\right\} \mathbb{E}_{R}\left[\mu_{jk}\tau_{rk}\right]\nonumber \\
 & =\sum_{r}\mathbb{E}_{R}\left[\mu_{jk}\mu_{jl}\rvert s_{j}=r\right]P_{R}\left(s_{j}=r\right)\tau_{rl}\label{eq:sector-effect}
\end{align}
where the $\mu_{jk}\tau_{rk}$ terms sum to zero from the normalization
of the sector level discounts in Equation \ref{eq:interchange-fee-model}.

This formula says that the sector weighted discounts are transmitted
sector-by-sector according to how much consumers overlap in each sector.
Intuitively, if debit card and credit card consumers overlap more
in groceries than in travel, then a discount for groceries and a higher
fee for travel benefits debit card users at the cost of credit card
users.

\paragraph{Negotiated Rates}

We model bargaining by simultaneously removing the bargaining parameter
$\delta_{l}$ while increasing the baseline fee $\overline{\tau}_{l}$
so as to maintain the same overall level of fees on card $l$. In
effect, we model a removal of price discrimination along the firm
size dimension. The dollar effect of eliminating $\delta_{l}$ while
increasing $\overline{\tau}_{l}$ to compensate is:
\begin{align}
 & -\delta_{l}\times\left(\frac{\partial V_{k}}{\partial\delta_{l}}-\frac{\partial V_{k}}{\partial\overline{\tau}_{l}}\times\frac{\mathbb{E}_{R}\left[\mu_{jl}b_{j}\right]}{\mathbb{E}_{R}\left[\mu_{jl}\right]}\right)\nonumber \\
= & \delta_{l}P_{R}\left(b_{j}\right)\left(\mathbb{E}_{R}\left[\mu_{jk}\mu_{jl}\rvert b_{j}=1\right]-\mathbb{E}_{R}\left[\mu_{jk}\mu_{jl}\right]\times\frac{\mathbb{E}_{R}\left[\mu_{jl}\rvert b_{j}=1\right]}{\mathbb{E}_{R}\left[\mu_{jl}\right]}\right)\label{eq:bargain-effect}
\end{align}

This formula predicts that debit card users benefit from the discounts
that large merchants receive on credit card interchange provided that
debit and credit card users overlap more at the merchants that receive
discounts than at the merchants that do not. To see this, consider
the case when bargaining and non-bargaining firms have the same average
use of payment method $l$, so that $\frac{\mathbb{E}_{R}\left[\mu_{jl}\rvert b_{j}=1\right]}{\mathbb{E}_{R}\left[\mu_{jl}\right]}=1$.
Then if $\delta_{l}<0$, consumers on card $k$ benefit from the discount
(i.e. are hurt if the discount is removed) provided that $\mathbb{E}_{R}\left[\mu_{jk}\mu_{jl}\rvert b_{j}=1\right]>\mathbb{E}_{R}\left[\mu_{jk}\mu_{jl}\rvert b_{j}=0\right]$.
This condition holds if $k$ users overlap more with $l$ users at
the large stores with discounted interchange than at the small stores
that do not have discounted interchange. Because large merchants have
a consumer base that covers consumers with different income levels,
this condition will generally be satisfied, highlighting a progressive
role for negotiated rates.

\subsubsection{Comments on the Sufficient Statistic}

The sufficient statistic formula \ref{eq:sufficient-stat-formula}
shows that the redistributive effects of interchange fees can be summarized
by the distribution of of payment shares $\mu$ across merchants provided
two conditions hold: full merchant passthrough to retail prices and
full issuer passthrough to rewards.

The sufficient statistic is valid even when consumers are allowed
to re-optimize their consumption in response to retail prices or consumer
rewards. One intuitive way of evaluating the effects of reducing interchange
fees would be to conduct an accounting exercise in which one adjusts
retail prices at each store according to the weighted change in interchange
fees, and then take the expenditure share-weighted sum of these price
changes to compute welfare. The formula in equation \ref{eq:sufficient-stat-formula}
provides a numerically convenient way of doing this exercise. One
concern with such an exercise is that it omits the result of consumers
responding to the change in retail prices by changing their shopping
behavior. The envelope theorem and Roy's identity highlight that such
concerns are second order in the change in interchange fees $\tau$.

The validity of the sufficient statistic does not depend on any particular
value of merchant demand elasticity. The essential ingredient comes
from the passthrough of interchange fees to retail prices. The best
empirical evidence that suggests interchange fees should be largely
passed through comes from \citet{Amiti.etal2019}, who show that sector-level
cost shocks are fully passed through, whereas idiosyncratic cost shocks
are only passed through incompletely at large firms. Given that interchange
fees affect the cost of all merchants, the usual evidence of incomplete
passthrough that typically applies to one firm at a time does not
apply to the effects of interchange.

Although the framework treats rewards as being paid out as a lump-sum,
this is not essential to the argument. If rewards were paid out for
specific merchants or sectors, then taking away rewards would change
consumers' allocation of spending across merchants. But by Roy's identity,
the re-optimization across merchants is second-order for welfare.
Only the income effect while holding fixed consumers' consumption
decisions is directly relevant.

\subsection{Counterfactual: Quantifying Redistribution}

We utilize our general formulas to examine the welfare effects of
setting interchange fees to zero, enabling us to understand the distributional
impact of interchange fees across different consumer payment types.
We apply our framework to a cross-section of sales and interchange
fees by card type in 2022. To implement our sufficient statistics
framework, we simply need to know, at the merchant level, equilibrium
payment shares and the corresponding interchange fees, both of which
are directly observable in our data.

\subsubsection{Eliminating Interchange Fees and Card Rewards\label{subsec:eliminate-interchange}}

Eliminating interchange fees and card rewards has two equal and opposite
effects on consumers as a whole. On one hand, when interchange fees
are set to zero, consumers no longer receive rewards when they pay
with debit or credit cards. On the other hand, eliminating interchange
fees reduces merchants\textquoteright{} transaction costs, which leads
to lower prices, benefiting all consumers. Although these effects
offset each other in the aggregate, they have large distributional
effects. Consumers who use high-reward, and therefore high-interchange
credit cards, are worse off, whereas consumers who use low-fee payment
methods like cash and regulated debit cards are better off. Moreover,
the consumers who use low-fee payment methods at stores that primarily
serve other consumers that use high-fee payment methods experience
the largest benefits.

We compute the consumer welfare effects of setting interchange fees
to zero after accounting for consumer sorting and merchant fee heterogeneity
in Table \ref{tab:redist}. We calculate effects measured both as
a percentage of each consumer group's baseline expenditures as well
as in terms of overall dollars when the total expenditures across
all consumers are scaled up to around $\$12$ trillion. On average,
cash-using consumers would be better off by \scalarinput{combined_cash_final}
basis points of baseline consumption. Across all cash users, this
amounts to approximately \$\scalarinput{combined_dollar_cash_final}
billion in additional annual consumption. Cash users benefit purely
because the decline in interchange fees lowers effective prices, and
they are not hurt by the loss of rewards, since they never received
them in the first place.

Even though much of the literature focuses on transfers between card
and cash users, regulated debit card users also benefit substantially
from capping interchange fees. For this group, average consumer surplus
would increase by \scalarinput{combined_regulated_debit_final} basis
points, corresponding to approximately \$\scalarinput{combined_dollar_regulated_debit_final}
billion in additional annual consumption. Although these users lose
out on modest card rewards, the associated price reductions more than
offset the loss. This is because rewards tied to regulated debit cards
are relatively small, especially compared to those from credit cards,
and credit card usage is more prevalent. As a result, setting interchange
fees to zero yields a significant decline in average prices relative
to the value of regulated debit card rewards. These findings indicate
that it is not only cash users who subsidize credit card rewards.
Regulated debit card users, who face lower fees and receive lower
rewards, also contribute to this subsidy.

Whereas cash and regulated debit card users benefit from reductions
in interchange fees, both basic and premium credit card users are
hurt by the reduction in interchange fees. This is because they are
disproportionately hurt by the reduction in rewards while receiving
only modest benefits from the reduction in retail prices. On net,
basic and premium credit card users lose \absinput{combined_basic_credit_final}
and \absinput{combined_premium_credit_final} bps of consumption,
respectively, amounting to \$\absinput{combined_dollar_basic_credit_final}
and \$\absinput{combined_dollar_premium_credit_final} billion, respectively.

It is helpful to put these changes in consumption in perspective.
One natural benchmark is the sales tax, which serves as a tax on consumption.
As of 2024, the average state sales tax rate in the U.S. was 6\%.\footnote{https://taxfoundation.org/data/all/state/sales-tax-revenue-reliance-breadth/}
Our results suggest that the cross-subsidization between cash and
credit card users is equivalent to cash users effectively paying a
\scalarinput{effective_tax} percent higher sales tax than premium
credit card users. The numbers also imply that credit card users receive
around \$\absinput{combined_dollar_headline} billion in aggregate
from cash and debit card users. These annual transfers are large relative
to the outlays of many government social insurance programs, such
as SNAP benefits (\$120 billion), EITC (\$57 billion), and regular
unemployment insurance (\$40 billion). The transfers here are also
large relative to other transfers in consumer financial markets, such
as the inter-regional transfer due to the GSE's lack of risk-based
pricing (\$15 billion), consumers' losses from high-fee index funds
(\$20 billion), and consumer losses from shrouded credit card fees
(\$10 billion) \citep{agarwal2015regulating,Hurst.etal2016}.

These results also shed light on who receives rewards and who ``pays''
for those rewards. The orange squares in Figure \ref{fig:scatter_interchange}
display the share of total card rewards received by regulated and
exempt debit card users and basic and premium credit card users. The
share of rewards received reflects both card usage and associated
reward levels. The results indicate that premium credit card users
capture the largest share of total rewards at \scalarinput{premium_credit_rewards}\%.
In comparison, basic credit card users receive \scalarinput{basic_credit_rewards}\%,
and exempt and regulated debit card users receive \scalarinput{regulated_debit_rewards}\%
and \scalarinput{exempt_debit_rewards}\%, respectively.

The green circles in Figure \ref{fig:scatter_interchange} display
the share of interchange costs borne by cash, debit, and credit card
users. The results indicate that cash users pay for \scalarinput{cash_fees_heterogeneous_merchants}\%
of interchange fees and thus implicitly pay for \scalarinput{cash_fees_heterogeneous_merchants}\%
of card rewards, representing a transfer from cash users to card users.
In contrast, regulated debit and premium credit card users pay for
\scalarinput{regulated_debit_fees_heterogeneous_merchants}\% and
\scalarinput{premium_credit_fees_heterogeneous_merchants}\%, respectively,
of rewards programs.

Comparing the vertical distance between the share of rewards received
and the share of interchange fees paid highlights the extent of cross-subsidization.
For example, premium credit card users receive \scalarinput{premium_credit_rewards}\%
of rewards but pay only \scalarinput{premium_credit_fees_heterogeneous_merchants}\%
of total interchange fees. This \scalarinput{premium_diff} percentage
point gap represents a transfer from consumers using other payment
methods to premium credit card users. Conversely, regulated debit
users receive \scalarinput{regulated_debit_rewards}\% of rewards
but pay \scalarinput{regulated_debit_fees_heterogeneous_merchants}\%
of interchange fees. This implies that regulated debit cards substantially
subsidize the rewards programs associated with more expensive payment
methods such as credit cards. While the literature and policymakers
have often emphasized that cash users subsidize rewards programs,
our findings suggest that debit users, particularly regulated debit
card users, play a quantitatively important role in subsidizing credit
card rewards.

\subsubsection{Implications of Merchant Heterogeneity}

Our sufficient statistics framework shows that the sorting of consumers
with different payment preferences across merchants, along with fee
heterogeneity across merchants shapes the impact of interchange fees
on consumers. We find that accounting for merchant heterogeneity has
important implications for two reasons. First, as documented in Section \ref{sec:costs},
we find evidence of consumer sorting: cash, debit card, and credit
card users tend to shop at different merchants, which limits cross-subsidization.
Second, the merchants frequented by all consumer types, where the
potential for redistribution is the greatest, tend to pay the lowest
interchange fees, which inherently restricts redistribution.

To evaluate the relative importance of these channels, it is useful
to compare our model's results to a counterfactual model that instead
assumes all merchants serve the same mix of consumers, and that all
cards of a given card type charge all merchants the same fee. The
second row of Table \ref{tab:redist} shows the results under these
two assumptions about merchant homogeneity. If all merchants were
homogeneous, the gains to cash users and regulated debit card users
from setting interchange fees to zero would be substantially higher.
Consumer surplus among cash users would increase by \scalarinput{combined_cash_naive}
basis points (approximately \$\scalarinput{combined_dollar_cash_naive}
billion). Similarly, surplus among premium credit card users would
decrease by \scalarinput{combined_premium_credit_naive} basis points
(roughly \$\scalarinput{combined_dollar_cash_naive} billion). In
a world of homogenous merchants, the losses for premium card users
would be \$\absinput{combined_dollar_premium_credit_overstate} larger
, and the gains for cash users would be \$\absinput{combined_dollar_cash_overstate}
billion larger.

Consumer sorting explains most of the difference between our baseline
results and the homogeneous-merchant counterfactual. Specifically,
merchant heterogeneity, arising from both consumer sorting and interchange
fee variation, reduces credit card (basic and premium) surplus by
roughly \$9 billion relative to the homogeneous world. Of this, about
\$8 billion is explained by consumer sorting: because cardholders
disproportionately shop at card-accepting merchants, they effectively
fund a large share of their own rewards through higher prices. The
remaining \$1 billion comes from interchange fee heterogeneity: merchants
where cash, debit, and credit card users all shop---and where redistribution
could be greatest---tend to pay lower interchange fees, further limiting
redistribution.

Merchant heterogeneity, therefore, has important implications for
who pays for rewards. The purple triangles in Figure \ref{fig:scatter_interchange}
show the share of card interchange fees paid by each user group if
all merchants were homogeneous. This figure illustrates that cross-subsidization
would be substantially larger in a setting without merchant heterogeneity.
For example, premium credit card users would still receive \scalarinput{premium_credit_rewards}\%
of rewards, but they would pay for only \scalarinput{premium_credit_fees_homogenous_merchants}\%
of them. In other words, the degree of cross-subsidization for credit
card rewards would be much greater if there were no consumer sorting
across merchants (i.e., if all merchants were homogeneous).

The results indicate that premium credit card users would effectively
receive \scalarinput{premium_share_paid_homog}\% (=(\scalarinput{premium_credit_rewards}\%-\scalarinput{premium_credit_fees_homogenous_merchants}\%)/\scalarinput{premium_credit_rewards}\%)
of their rewards for free if merchants were homogeneous. However,
due to consumer sorting across merchants by payment method, premium
credit card users currently receive only \scalarinput{premium_share_paid_final}\%
(=(\scalarinput{premium_credit_rewards}\%-\scalarinput{premium_credit_fees_heterogeneous_merchants}\%)/\scalarinput{premium_credit_rewards}\%)
of their rewards for free. Thus, ignoring merchant heterogeneity overstates
the degree of cross-subsidization by around one-half. This highlights
the crucial role that consumer sorting and variation in interchange
rates play in reducing the extent of redistribution across payment
types.

\subsubsection{Redistribution by Income}

We also examine how the redistributive effects of interchange fees
and rewards programs vary across household income levels. In Section
\ref{subsec:eliminate-interchange}, we showed how much consumers
with different payment preferences (e.g., debit or credit) gain or
lose from capping interchange fees, expressed as a share of their
baseline consumption. To translate these results into distributional
effects by income, we combine them with the evidence in Figure \ref{fig:payment_income}
on how payment preferences vary with income. At each income level,
we calculate the net effect as a weighted average of the consumption
impacts across payment types, using as weights the market shares of
each payment method among households at that income level. We then
map these effects to dollar transfers by assuming consumers of all
income levels have the same proportion of purchases to income, and
by scaling up total purchases to $\$12$ trillion.

We find that capping interchange would generate a large progressive
transfer from households with incomes greater than $\$150,000$ to
those with incomes below that level. Figure \ref{fig:redistribution_interchange}
shows that the highest income households lose around $10$ basis points
of consumption because they have the highest levels of credit card
use. As a result, they would be hurt by the decline in rewards more
than they would benefit from lower retail prices. In contrast, the
effect is reversed for lower income households who are more likely
to use cash and debit cards. In dollar terms, this is a transfer of
around $\$200$ per year from households with incomes above $\$150,000$
per year to fund transfers of around $\$30$ to households with incomes
below that cutoff. Cumulatively, the high income households lose $\$4$
billion from a cap on interchange fees, whereas the low-income households
gain around hte same amount.

\subsubsection{Effects of the Durbin Amendment}

Our framework also sheds light on the redistributive impact of the
Durbin Amendment. The Durbin Amendment was designed to limit interchange
fees on debit card transactions. Table \ref{tab:durbin} presents
a counterfactual in which all interchange fees are set according to
unregulated schedules (i.e., no regulated debit). The results indicate,
unsurprisingly, that regulated debit card users would substantially
benefit from deregulation, with their surplus increasing by \scalarinput{durbin_regulated_debit}
basis points of consumption (\$\scalarinput{durbin_regulated_debit_dollar}
billion). More interesting, however, is identifying which users benefit
from the Durbin Amendment, or equivalently, which users would be harmed
by its repeal. Collectively, credit card users benefit the most, with
their surplus increasing by \$4.4 billion. Thus, regulating interchange
fees on debit cards primarily benefited credit card users.

We also examine the redistributive impact of the Durbin Amendment
across household income levels. Figure \ref{fig:redistribution_durbin}
illustrates the effects of the Durbin Amendment. The figure shows
that low-income consumers were hurt the most by the Durbin Amendment,
at the expense of higher-income consumers. This pattern arises because
higher-income consumers are more likely to use credit cards, and thus
benefit from the Durbin Amendment, while lower-income consumers are
more likely to rely on regulated debit cards. The Durbin Amendment
reduced merchant costs and perks for debit card users, such as free
checking accounts. In equilibrium, high income consumers benefit from
lower retail prices, whereas lower income consumers were hurt by the
lost checking account perks \citep{Mukharlyamov.Sarin2025}. The results
indicate that the Durbin Amendment reduced the surplus of the lowest-income
consumers, particularly those with incomes below \$30,000.

\subsubsection{Who Pays for Premiumization?}

Lastly, we examine the welfare effects of the rise in premium credit
cards. As documented in Section \ref{sec:costs} and Figure \ref{fig:premium_credit_ts},
the share of premium credit cards increased from 15\% in 2006 to over
60\% in 2022. Since the average interchange rate on premium cards
is 2.34\%, compared to 1.64\% for basic cards, this shift has raised
total credit card interchange fees by roughly 20\%.

Table \ref{tab:premiumization} reports the redistributive effects
of the rise in premium credit cards compared to a benchmark where
there are no premium credit cards. Figure \ref{fig:redistribution_premiumization}
displays the corresponding redistributive impact of increased premiumization
across household income levels. Unsurprisingly, premium cardholders
benefit from the increase in premiumization: their surplus increases
by \scalarinput{premiumization_premium_credit}  basis points of consumption
(about \$\scalarinput{premiumization_premium_credit_dollar}  billion).
By contrast, cash users, exempt debit users, regulated debit users,
and basic cardholders all are hurt by the rise in premiumization.
Interestingly, the largest costs accrue not to cash users, but to
regulated debit and basic credit card users. These groups shop more
frequently at the same merchants as premium cardholders, and therefore
bear a greater share of the burden from the higher interchange fees
associated with premium cards.

This finding underscores that it is not only cash users who subsidize
credit card rewards; low-reward cardholders also subsidize the rewards
of high-reward cardholders. Moreover, because all consumers face the
same prices regardless of payment method, consumers have incentives
to adopt the most expensive card offering the most generous rewards.
Consequently, costly payment methods, such as premium credit cards,
tend to crowd out cheaper ones, such as basic cards, helping to explain
the rise of premiumization.

\section{Conclusion\label{sec:conclusion}}

We show that interchange fees are not just a technical feature of
the payments industry but a pervasive force shaping retail activity
and redistributing consumption among consumers. First, using merchant-level
data, we document substantial heterogeneity in the composition and
cost of payments across stores. This heterogeneity has critical implications
for merchants\textquoteright{} costs and redistribution. We also show,
using the Durbin Amendment as a natural experiment, that interchange
fees matter for real outcomes: a one--percentage point decline in
interchange rates causally increases card sales by about 7 percent
over two years.

Finally, through a sufficient-statistics framework, we quantify the
overall impact of interchange pricing, finding that interchange fees
redistribute approximately \$30 billion per year from cash and debit
users to credit card users, with consumer sorting attenuating these
transfers by about 25 percent. We also find that the heterogeneity
in fees across sectors and the firm-size distribution are progressive,
redistributing around \$ \scalarinput{low_income_firm_redist} billion
per year to cash and exempt debit card users. In addition, we document
that two recent shifts in the industry, the Durbin Amendment and the
rise of premium credit cards, have had regressive impacts. Surprisingly,
and likely unintentionally, the Durbin Amendment primarily benefited
credit card users.These results highlight that interchange policies
affect not only banks and card networks but also prices, market competition,
and the distribution of economic surplus across consumers.

\clearpage{}

\begin{singlespace}
\begin{flushleft}
{\small\bibliographystyle{chicago}
\bibliography{payment_methods,payments_fiserv_lulu}
}{\small\par}
\par\end{flushleft}
\end{singlespace}

\clearpage{}

\section*{Figures}

\begin{figure}[H]
\caption{Illustration of payment flows in a payment network\label{fig:payment-flow}}
\small

\[
\xymatrix{ & *+++++[F:<5pt>]{\text{Network}}\\*
+++++[F:<5pt>]{\text{Acquirer}}\ar@(u,l)[ru]_{0.07}^{\text{Assessment\ }}\ar[rr]_{1.90}^{\text{Interchange}} &  & *+++++[F:<5pt>]{\text{Issuer}}\ar@(u,r)[lu]_{\text{\ Assessment}}^{0.07}\ar[d]_{\text{Rewards}}^{1.40}\\*
+++++[F:<5pt>]{\text{Merchant}}\ar[u]_{2.25}^{\text{Merchant Discount}} &  & *+++++[F:<5pt>]{\text{Consumer}}\ar[ll]_{\text{Purchase}}^{100}
}
\]

\caption*{\emph{Notes: }} Figure \ref{fig:payment-flow} shows average
fee levels for a credit card transaction in the US. Data on merchant
fees comes from \citep{Nilson2023}. \citet{Drechsler.etal2025} calculate
average interchange fees and rewards from regulatory data on large
banks in the US. The average network fee comes from fees disclosed
by acquirers as well as the annual financial statements from Visa
\citep{Visa2020,Helcim2021}.

\end{figure}

\clearpage{}
\begin{figure}[H]
  \centering
  \caption{Data Coverage and Representativeness}
  \label{fig:data-representative}

  % Row 1
  \begin{subfigure}[t]{0.48\textwidth}
    \centering
    \includegraphics[width=\linewidth]{../output/graphs/dcpc_settlement_compare.pdf}
    \caption{Sector Composition}
    \label{fig:data-representative-sector}
  \end{subfigure}\hfill
  \begin{subfigure}[t]{0.48\textwidth}
    \centering
    \includegraphics[width=\linewidth]{../output/graphs/state_scatter.pdf}
    \caption{Geographic Composition}
    \label{fig:data-representative-geo}
  \end{subfigure}

  % Row 2
  \vspace{0.6em}
  \begin{subfigure}[t]{0.65\textwidth}
    \centering
    \includegraphics[width=\linewidth]{../output/graphs/firm_size_distn_irs_compare.pdf}
    \caption{Firm Size}
    \label{fig:data-representative-size}
  \end{subfigure}

  \captionsetup{font=small}
  \caption*{\textit{Notes:} Figure \ref{fig:data-representative-sector} shows the share (by value)
  of card transactions across sectors according to different datasets. DCPC refers to data from the
  Atlanta Fed's Diary of Consumer Payment Choice. Settlement refers to the main Fiserv card settlement
  data that covers both large and small merchants.

  Figure \ref{fig:data-representative-geo} shows the GDP share of each state versus the share (by value)
  of card transactions across states in the Fiserv settlement data.

  Figure \ref{fig:data-representative-size} shows the share (by value) of transactions by merchants with
  different levels of revenue. The IRS measure uses the Statistics on Income to show the share of gross
  receipts in each firm-size category. The Fiserv measure uses the settlement data to show the share of
  card spending in each firm-size category.}
\end{figure}\clearpage{} 
\begin{figure}[H]
  \centering
  \caption{Heterogeneity in the Payment Mix Across Merchants}
  \label{fig:pmt_merchant}

  % Row 1
  \begin{subfigure}[t]{0.48\textwidth}
    \centering
    \includegraphics[width=\linewidth]{../output/graphs/clover_cash_condl_share_hist.pdf}
    \caption{Cash Share}
    \label{fig:cash_share}
  \end{subfigure}\hfill
  \begin{subfigure}[t]{0.48\textwidth}
    \centering
    \includegraphics[width=\linewidth]{../output/graphs/clover_cash_share_ts.pdf}
    \caption{Cash Share over Time}
    \label{fig:cash_share_ts}
  \end{subfigure}

  % Row 2
  \vspace{0.6em}
  \begin{subfigure}[t]{0.48\textwidth}
    \centering
    \includegraphics[width=\linewidth]{../output/graphs/merchant_cs_credit_share.pdf}
    \caption{Credit Card Share of Card Payments}
    \label{fig:credit_share}
  \end{subfigure}\hfill
  \begin{subfigure}[t]{0.48\textwidth}
    \centering
    \includegraphics[width=\linewidth]{../output/graphs/credit_share_over_time_annual.pdf}
    \caption{Credit Card Share over Time}
    \label{fig:credit_share_ts}
  \end{subfigure}

  \caption*{\small \textit{Notes:} Figure \ref{fig:cash_share} shows the distribution of cash and check
  share among merchants in the Clover data, conditional on having at least 2\% of payments made via
  cash or check. Around two-thirds of merchant-year observations do not feature any cash transactions.
  Observations are at the merchant-year level and are equally weighted over the period 2019--2022.\\
  Figure \ref{fig:cash_share_ts} displays the corresponding dollar-weighted average cash share over time,
  computed at the annual level.\\
  Figure \ref{fig:credit_share} shows the distribution of the dollar share of card transactions on credit
  cards across merchants in 2022 with more than \sizecutoff\ in annual card sales. Observations are at the
  merchant-year level and are equally weighted.\\
  Figure \ref{fig:credit_share_ts} displays the corresponding dollar-weighted average credit card share of
  total card transactions over time. The time series excludes large firms in the merchandise and retail
  sector with more than \$1 billion in sales to prevent the entry and exit of large merchants from the
  sample from influencing trends in fees.}
\end{figure}\clearpage{}


\begin{figure}[H]
\centering
\caption{Heterogeneity in the Payment Mix Across Sectors} \label{fig:pmt_sectors}

\includegraphics[width=6in]{../output/graphs/payment_share_cash_sector.pdf}

\caption*{\small \textit{Notes:} Figure \ref{fig:pmt_sectors}  shows the dollar share of different types of card transactions across major merchant categories in the Clover data.  }

\end{figure}

\clearpage{}


\begin{figure}[H]
\centering
\caption{Variation in Composition of Card Payments Within Sectors} \label{fig:sector_densities}

\includegraphics[width=6in]{../output/graphs/merchant_cs_credit_share_by_sector.pdf}

\caption*{\small \textit{Notes:} Figure \ref{fig:sector_densities}  shows credit card transactions as a dollar share of total card transactions across firms in 2022.}

\end{figure}\clearpage{}
\begin{figure}[H]
  \centering
  \caption{Geographic Variation in the Payment Mix}
  \label{fig:pmt_map}

  \begin{subfigure}[t]{0.9\textwidth}
    \centering
    \includegraphics[width=\linewidth]{../output/graphs/map_cash.png}
    \caption{Cash Share}
    \label{fig:cash_map}
  \end{subfigure}

  \vspace{0.6em}

  \begin{subfigure}[t]{0.9\textwidth}
    \centering
    \includegraphics[width=\linewidth]{../output/graphs/map_credit.png}
    \caption{Credit Card Share of Card Payments}
    \label{fig:credit_map}
  \end{subfigure}

  \captionsetup{font=small}
  \caption*{\textit{Notes:} Figure \ref{fig:pmt_map} shows the geographic distribution of cash and credit
  card payments. Figure \ref{fig:cash_map} displays the share of all transactions conducted in cash,
  equally weighted across merchants in the county. Figure \ref{fig:credit_map} presents the share of total
  card transactions made using credit cards across U.S. counties, equally weighted across merchants in the
  county. Lighter colors in panels (a) and (b) indicate regions with higher shares of cash and credit card
  payments, respectively.}
\end{figure}

\clearpage{}
\begin{figure}[H]
  \centering
  \caption{Payment Choice and Income}
  \label{fig:payment_income}

    \includegraphics[width=0.9\linewidth]{../output_nonfiserv/graphs/payment_income.pdf}
    \caption{Preferences for Cash, Debit, and Credit by Income}

  \captionsetup{font=small}
  \caption*{\textit{Notes:} Figure \ref{fig:payment_income} shows data from MRI's Ultimate Survey of Americans on consumer's primary payment methods by income. Cash respondents are those who agree with the statement "I try to pay with cash, when possible". Debit card consumers are non-cash consumers who report more monthly spending on debit cards than credit cards. Within debit cards, small bank consumers are those who report having an account at a community bank or credit union. Within credit cards, premium credit consumers are those who report owning a Visa Signature, MC world elite, or American Express card.}
\end{figure}\clearpage{}

\begin{figure}[H]
\centering

\caption{Distribution of Avg. Transaction Fees (Excluding Cash)} \label{fig:interchange_dist_card}

\includegraphics[width=5in]{../output/graphs/merchant_cs_avg_interchange.pdf}

\caption*{\small \textit{Notes:} Figure \ref{fig:interchange_dist_card} shows the distribution of average interchange fees across merchants. The sample covers merchants in 2022 with more than \sizecutoff in card sales.}


\end{figure}\clearpage{}

\begin{figure}[H]
  \centering
  \caption{Interchange Fees by Instrument}
  \label{fig:interchange_dist}

  % Row 1
  \begin{subfigure}[t]{0.48\textwidth}
    \centering
    \includegraphics[width=\linewidth]{../output/graphs/merchant_cs_icg_distn_by_type.pdf}
    \caption{Distribution of Fees by Instrument}
    \label{fig:interchange_by_pmt}
  \end{subfigure}\hfill
  \begin{subfigure}[t]{0.48\textwidth}
    \centering
    \includegraphics[width=\linewidth]{../output/graphs/debit_credit_icg_rate_annual.pdf}
    \caption{Average Debit and Credit Fees over Time}
    \label{fig:interchange_ts}
  \end{subfigure}

  % Row 2
  \vspace{0.6em}
  \begin{subfigure}[t]{0.48\textwidth}
    \centering
    \includegraphics[width=\linewidth]{../output/graphs/debit_icg_annual.pdf}
    \caption{Debit Card Fees over Time}
    \label{fig:interchange_debit_ts}
  \end{subfigure}\hfill
  \begin{subfigure}[t]{0.48\textwidth}
    \centering
    \includegraphics[width=\linewidth]{../output/graphs/credit_icg_annual.pdf}
    \caption{Credit Card Fees over Time}
    \label{fig:interchange_credit_ts}
  \end{subfigure}

  \captionsetup{font=small}
  \caption*{\textit{Notes:} Figure \ref{fig:interchange_by_pmt} shows kernel density plots of the average
  interchange fee for different payment methods across merchants with more than \$100{,}000 in sales in 2022.
  The leftmost distribution is for regulated debit cards; the dotted red distribution is for exempt debit cards;
  the dashed green distribution is for basic credit cards; and the peaked orange distribution is for premium
  credit cards.\\
  Figure \ref{fig:interchange_ts} shows the evolution of dollar-weighted average credit- and debit-card
  interchange fees over time.\\
  Figure \ref{fig:interchange_debit_ts} shows how dollar-weighted debit-card interchange fees vary by debit
  card type over time. The \textit{Regulated} line reflects the lower interchange tier created after the 2011
  Durbin Amendment.\\
  Figure \ref{fig:interchange_credit_ts} shows how dollar-weighted credit-card interchange fees vary by credit
  card type over time; premium credit cards target transactors and earn higher interchange fees.\\
  All time series exclude large firms in the merchandise and retail sector with more than \$1~billion in
  sales to prevent entry and exit of large merchants from influencing fee trends.}
\end{figure}\clearpage{}


\begin{figure}[H]
  \centering
  \caption{Composition of Card Payments}
  \label{fig:card_type_ts}

  \begin{subfigure}[t]{0.7\textwidth}
    \centering
    \includegraphics[width=\linewidth]{../output/graphs/debit_reg_share_over_time_annual.pdf}
    \caption{Share of Debit that is Regulated}
    \label{fig:regulated_debit_ts}
  \end{subfigure}

  \vspace{0.6em}

  \begin{subfigure}[t]{0.7\textwidth}
    \centering
    \includegraphics[width=\linewidth]{../output/graphs/premium_share_over_time_annual.pdf}
    \caption{Share of Credit that is Premium}
    \label{fig:premium_credit_ts}
  \end{subfigure}

  \captionsetup{font=small}
  \caption*{\textit{Notes:} Figure \ref{fig:regulated_debit_ts} shows the dollar-weighted share over time
  of debit card transactions that receive capped interchange. Figure \ref{fig:premium_credit_ts} shows the
  dollar-weighted share over time of credit card transactions that receive premium interchange.
  Both time series exclude large firms in the merchandise and retail sector with more than \$1 billion in
  sales to prevent the entry and exit of large merchants from the sample from influencing trends in fees.}
\end{figure}

\clearpage{}

\begin{figure}[H]
  \centering
  \caption{Variation in Interchange Fees by Card Type, Merchant Sector, and Size}
  \label{fig:price-disc}

  % Row 1
  \begin{subfigure}[t]{0.48\textwidth}
    \centering
    \includegraphics[width=\linewidth]{../output/graphs/disc_card_type.pdf}
    \caption{Interchange Fees by Card Type}
    \label{fig:price-disc-card}
  \end{subfigure}\hfill
  \begin{subfigure}[t]{0.48\textwidth}
    \centering
    \includegraphics[width=\linewidth]{../output/graphs/disc_sector.pdf}
    \caption{Interchange Fees by Sector}
    \label{fig:price-disc-merch}
  \end{subfigure}

  % Row 2
  \vspace{0.6em}
  \begin{subfigure}[t]{0.65\textwidth}
    \centering
    \includegraphics[width=\linewidth]{../output/graphs/disc_negotiated.pdf}
    \caption{Interchange Fees by Merchant Size}
    \label{fig:price-disc-size}
  \end{subfigure}

  \captionsetup{font=small}
  \caption*{\textit{Notes:} Figure \ref{fig:price-disc-card} shows dollar-weighted interchange fees for
  different card types for merchants with more than \sizecutoff\ in sales, excluding merchants that receive
  volume discounts. Premium credit cards are defined as transactions on the Visa Signature and Visa Infinite
  (and Mastercard equivalents) interchange tiers. Regulated debit cards refer to debit cards issued by banks
  with more than \$10~billion in assets; exempt debit cards are those issued by smaller banks. Interchange
  fees are calculated as total interchange fees paid divided by total card volume within the category.
  Figure \ref{fig:price-disc-merch} shows average interchange fees across sectors for credit cards and exempt
  debit cards for firms that do not receive volume discounts. Figure \ref{fig:price-disc-size} provides
  evidence of volume discounts by comparing average interchange fees by merchant size in the grocery and
  merchandise categories. When calculating size, we use the merchant's total card sales across all
  establishments.}
\end{figure}

\clearpage{}


\begin{figure}[H]
\centering
\caption{Variation in Merchant Interchange Fees Explained by Different Factors} \label{fig:interchange_decomp}

\includegraphics[width=6in]{../output/graphs/merchant_cs_icg_explained.pdf}

\caption*{\small \textit{Notes:} Figure \ref{fig:interchange_decomp} shows the distribution of average interchange fees across merchants in solid blue, and then shows how much of the variation in interchange fees can be explained by different factors. }


\end{figure}

\clearpage{} 
\begin{figure}[H]
  \centering
  \caption{Explaining the Heterogeneous Impact of the Durbin Amendment Across Restaurants}
  \label{fig:durbin-het-inputs}

  \begin{subfigure}[t]{0.7\textwidth}
    \centering
    \includegraphics[width=\linewidth]{../output/graphs/durbin_ticket_size_effect.pdf}
    \caption{Variation by Ticket Size}
    \label{fig:durbin-het-ticket}
  \end{subfigure}

  \vspace{0.6em}

  \begin{subfigure}[t]{0.7\textwidth}
    \centering
    \includegraphics[width=\linewidth]{../output/graphs/durbin_regulated_share_effect_zip.pdf}
    \caption{Variation by Share of Regulated Debit Cards}
    \label{fig:durbin-het-regulated}
  \end{subfigure}

  \captionsetup{font=small}
  \caption*{\textit{Notes:} Figure \ref{fig:durbin-het-inputs} illustrates the heterogeneous impact of the
  Durbin Amendment across merchants in the restaurant sector (MCC codes 5812 and 5814) and across regions.
  Panel \ref{fig:durbin-het-ticket} shows a binscatter of pre-Durbin average ticket size against pre- and
  post-Durbin debit interchange rates. Ticket size is calculated as the total dollar value of debit card
  transactions divided by the total count of debit card transactions in the October 2010--September 2011
  window preceding the Durbin Amendment; observations are at the merchant-year level. Panel
  \ref{fig:durbin-het-regulated} shows a binscatter of the ZIP-code-level post-Durbin share of debit cards
  that earn regulated interchange against pre- and post-Durbin debit interchange rates.}
\end{figure}\clearpage{}

\begin{figure}[H]   
\caption{Who Pays for and Receives Rewards}    
\label{fig:scatter_interchange}
\centering   
\small
\includegraphics[width=5in]{../output/graphs/shares_scatter_by_method.pdf}
\caption*{\footnotesize   Figure \ref{fig:scatter_interchange} displays the share of interchange fees paid and rewards received by consumers of different payment methods. The figure shows how the burden of interchange fees changes after modeling merchant heterogeneity in both payment composition and fees. The figure is calculated under the assumption that all merchants passthrough interchange fees into prices, dollar for dollar. }
\end{figure}

\clearpage{}
\begin{figure}[H]
  \centering
  \caption{Redistribution Across Household Income}
  \label{fig:redistribution}

  \begin{subfigure}[t]{0.9\textwidth}
    \centering
    \includegraphics[width=\linewidth]{../output_nonfiserv/graphs/cap_bps.pdf}
    \caption{Eliminating Interchange Fees}
    \label{fig:redistribution_interchange}
  \end{subfigure}


  \begin{subfigure}[t]{0.45\textwidth}
    \centering
    \includegraphics[width=\linewidth]{../output_nonfiserv/graphs/durbin_bps.pdf}
    \caption{Durbin Amendment}
    \label{fig:redistribution_durbin}
  \end{subfigure}
  \begin{subfigure}[t]{0.45\textwidth}
    \centering
    \includegraphics[width=\linewidth]{../output_nonfiserv/graphs/premiumization_bps.pdf}
    \caption{Rise of Premium Credit Cards}
    \label{fig:redistribution_premiumization}
  \end{subfigure}

  \captionsetup{font=small}
\caption*{\textit{Notes:} Figure \ref{fig:redistribution} illustrates the impact of policy changes on household surplus by income. Panel (a) shows the change in household surplus if interchange fees and card reward programs were banned. Panel (b) shows the change in household surplus if the Durbin Amendment were repealed. Panel (c) shows the change in household surplus if premium credit cards were banned.}

\end{figure}
\begin{center}
\clearpage{}
\par\end{center}

\section*{Tables}


\begin{table}[H]
  \centering
  \caption{Summary of Fiserv Data}
  \label{tab:fiserv-summary}

  \begin{subtable}[t]{0.95\textwidth}
    \centering
    \caption{Settlement Data Cross Section from 2022}
    \label{tab:settlement-summary}
    \small
    \input{../output/tables/settlement_summary.tex}
  \end{subtable}

  \vspace{0.8em}

  \begin{subtable}[t]{0.95\textwidth}
    \centering
    \caption{Restaurant Settlement Panel 2010--2013}
    \label{tab:durbin-summary}
    \small
    \input{../output/tables/durbin_summary.tex}
  \end{subtable}

  \vspace{0.8em}

  \begin{subtable}[t]{0.95\textwidth}
    \centering
    \caption{Summary of Clover Data}
    \label{tab:clover-summary}
    \small
    \input{../output/tables/clover_summary.tex}
  \end{subtable}

  \captionsetup{font=small}
  \caption*{\emph{Notes:} Panel (a) shows summary statistics from the 2022 cross-section at the
  merchant-sector level. The payment-method variables refer to the share of card sales at the merchant
  on each payment method. Regulated debit cards are issued by large banks. Exempt debit cards are issued
  by small banks. Premium credit refers to high-interchange credit cards in the Visa Signature and Visa
  Infinite (and the Mastercard equivalent) interchange tiers. Average interchange is calculated as total
  interchange paid divided by total sales. Transaction size is the total dollar volume divided by the
  number of transactions. Panel (b) shows summary statistics at the merchant-year level from our
  2010--2013 panel of restaurant sales. Regulated debit statistics are computed for 2011 and onward.
  Interchange rates are calculated as total interchange within a category divided by total sales in that
  category. Panel (c) shows summary statistics of the Clover data from 2019--2022 at the merchant-year
  level. The cash and check share is defined as the dollar value of cash and check transactions divided
  by the total value of transactions for the merchant in that year.}
\end{table}
\begin{center}
\clearpage{}
\begin{table}[H]
\caption{Spatial and Firm Determinants of Payment Mix \label{tab:payment_type}}

\begin{centering}
\small{\input{../output/tables/geo_tables.tex}}
\par\end{centering}
\caption*{\emph{Notes: }Columns (1) to (3) report regressions of
the firm-level share of transactions on cash and check in the Clover
data on firm characteristics, sectors, and spatial factors. Columns
(4) to (6) report regressions using the settlement data of the firm-level
share of card transactions on credit cards on firm characteristics,
sectors, and spatial factors.}
\end{table}
\par\end{center}

\begin{center}
\clearpage{}
\begin{table}[H]
\caption{Summary Statistics on Interchange Rates by Card Type Across Firms
\label{tab:icg-summary}}

\centering

\input{../output/tables/settlement_summary_rates}

\caption*{\emph{Notes: }Interchange rates are defined as total net
interchange divided by total net sales for a given card type. Interchange
fees are set as missing if the given payment method does not have
any transactions. Merchants are equal weighted, which explains differences
with the aggregate fees in Figure \ref{fig:price-disc}.}
\end{table}
\par\end{center}

\begin{center}
\par\end{center}

\begin{center}
\clearpage{}
\par\end{center}

\begin{center}

\begin{table}[H]
  \centering
  \caption{Interchange rates from Visa's 2022 interchange publication.}
  \label{tab:visa-interchange-public}

  \begin{subtable}[t]{0.95\textwidth}
    \centering
    \caption{By Sector and Card Type}
    \label{tab:public-icg-sector}
    \small
    \begin{tabular}{>{\raggedright\arraybackslash}p{1.6in}cccc}
      \hline
      Card Category & Gas Station & Supermarket & Merchandise & Restaurant \\
      \hline
      Regulated Debit Card & $0.05\%+0.22$ & $0.05\%+0.22$ & $0.05\%+0.22$ & $0.05\%+0.22$ \\
      Exempt Debit Card    & $\min\{0.8\%+0.15,\$0.95\}$ & \$0.30 & $0.80\%+0.15$ & $1.19\%+0.10$ \\
      Basic Credit         & $\min\{1.15\%+0.25,\$1.10\}$ & $1.50\%+0.07$ & $1.51\%+0.10$ & $\max\{2.10\%,\$0.04\}$ \\
      Premium Credit       & $\min\{1.15\%+0.25,\$1.10\}$ & $2.00\%+0.07$ & $2.10\%+0.10$ & $\max\{2.60\%,\$0.04\}$ \\
      \hline
    \end{tabular}
  \end{subtable}

  \vspace{0.8em}

  \begin{subtable}[t]{0.6\textwidth}
    \centering
    \caption{By Firm Size}
    \label{tab:public-icg-firm-size}
    \small
    \begin{tabular}{>{\raggedright\arraybackslash}p{2.8in}c}
      \hline
      Supermarket Size Category & Premium Credit Rate \\
      \hline
      Tier 0 $\left(\geq \$19.6\ \text{Bn}\right)$ & $1.40\% + \$0.05$ \\
      Tier 1 $\left(\geq \$10.65\ \text{Bn}\right)$ & $1.55\% + \$0.05$ \\
      Tier 2 $\left(\geq \$4.2\ \text{Bn}\right)$ & $1.65\% + \$0.05$ \\
      Tier 3 $\left(\geq \$0.95\ \text{Bn}\right)$ & $1.75\% + \$0.05$ \\
      All Other (Smallest) & $2.00\% + \$0.07$ \\
      \hline
    \end{tabular}
  \end{subtable}

  \captionsetup{font=small}
  \caption*{\emph{Notes:} Percents denote the variable part of the interchange fee; the remainder is the
  fixed fee. Thus $0.05\%+0.22$ indicates a fee of $0.05\%$ of transaction value plus \$0.22. Premium credit
  corresponds to Visa Infinite Spend Qualified Cards. Basic credit corresponds to ``All Other Products.''
  Regulated debit cards are those issued by large banks subject to the Durbin Amendment's caps on
  interchange fees.}
\end{table}\clearpage{}
\par\end{center}

\begin{center}
\begin{table}[!h]
\caption{The Effects of Changes in Interchange Fees on Sales and Prices\label{tab:durbin-iv}}

\begin{centering}
{\tiny\input{../output/tables/durbin_iv}}
\par\end{centering}
\raggedright{}\caption*{\emph{Notes: }Table \ref{tab:durbin-iv}
shows the results of regressing changes in log sales and changes in
log ticket size on the change in average interchange rates paid at
restaurants in the 1 year and 2 year windows after the implementation
the Durbin Amendment. Log ticket size is measured as the log of the
ratio of the dollar value of sales to the number of transactions.
The pre-debit share is the share of debit card transactions at the
restaurant from October 2010 to September 2011. The Change in Fees
is defined as the change in the average per transaction debit interchange
fee for a regulated debit card at the store, as computed from the
average ticket size measured from October 2010 to September 2011 and
the public debit interchange schedule prior to the Durbin Amendment.}
\end{table}
\clearpage{}
\par\end{center}

\begin{center}
\begin{table}[H]
  \centering
  \caption{Effects of Capping Interchange Fees}
  \label{tab:redist}

  \begin{subtable}[t]{0.95\textwidth}
    \centering
    \caption{As Basis Points of Baseline Consumption}
    \label{tab:redist-bps}
    \small
    \begin{tabular}{>{\raggedright\arraybackslash}p{0.1cm} >{\raggedright\arraybackslash}p{1in}
                    >{\centering\arraybackslash}p{0.7in} >{\centering\arraybackslash}p{0.7in}
                    >{\centering\arraybackslash}p{0.7in} >{\centering\arraybackslash}p{0.7in}
                    >{\centering\arraybackslash}p{0.7in}}
      \hline
       &  & \multicolumn{5}{c}{Consumer Group} \\
      \cline{3-7}
      \multicolumn{2}{c}{Mechanism} & Cash & Regulated Debit & Exempt Debit & Basic Credit & Premium Credit \\
      \hline
      \multicolumn{2}{>{\raggedright}p{1in}}{Full Model}
        & \scalarinput{combined_cash_final}
        & \scalarinput{combined_regulated_debit_final}
        & \scalarinput{combined_unregulated_debit_final}
        & \scalarinput{combined_basic_credit_final}
        & \scalarinput{combined_premium_credit_final} \\
      \multicolumn{2}{>{\raggedright}p{1in}}{Homogeneous}
        & \scalarinput{combined_cash_naive}
        & \scalarinput{combined_regulated_debit_naive}
        & \scalarinput{combined_unregulated_debit_naive}
        & \scalarinput{combined_basic_credit_naive}
        & \scalarinput{combined_premium_credit_naive} \\
      \hline
      \multicolumn{2}{>{\raggedright}p{0.9in}}{Difference}
        & \scalarinput{combined_cash_overstate}
        & \scalarinput{combined_regulated_debit_overstate}
        & \scalarinput{combined_unregulated_debit_overstate}
        & \scalarinput{combined_basic_credit_overstate}
        & \scalarinput{combined_premium_credit_overstate} \\
       & Consumer Sorting
        & \scalarinput{combined_cash_sorting}
        & \scalarinput{combined_regulated_debit_sorting}
        & \scalarinput{combined_unregulated_debit_sorting}
        & \scalarinput{combined_basic_credit_sorting}
        & \scalarinput{combined_premium_credit_sorting} \\
       & Merchant Fee Heterogeneity
        & \scalarinput{combined_cash_firm}
        & \scalarinput{combined_regulated_debit_firm}
        & \scalarinput{combined_unregulated_debit_firm}
        & \scalarinput{combined_basic_credit_firm}
        & \scalarinput{combined_premium_credit_firm} \\
      \hline
    \end{tabular}
  \end{subtable}

  \vspace{0.8em}

  \begin{subtable}[t]{0.95\textwidth}
    \centering
    \caption{In Dollar Terms (\$Bn)}
    \label{tab:redist-dollars}
    \small
    \begin{tabular}{>{\raggedright\arraybackslash}p{0.1cm} >{\raggedright\arraybackslash}p{1in}
                    >{\centering\arraybackslash}p{0.7in} >{\centering\arraybackslash}p{0.7in}
                    >{\centering\arraybackslash}p{0.7in} >{\centering\arraybackslash}p{0.7in}
                    >{\centering\arraybackslash}p{0.7in}}
      \hline
       &  & \multicolumn{5}{c}{Consumer Group} \\
      \cline{3-7}
      \multicolumn{2}{c}{Mechanism} & Cash & Regulated Debit & Exempt Debit & Basic Credit & Premium Credit \\
      \hline
      \multicolumn{2}{>{\raggedright}p{1in}}{Full Model}
        & \scalarinput{combined_dollar_cash_final}
        & \scalarinput{combined_dollar_regulated_debit_final}
        & \scalarinput{combined_dollar_unregulated_debit_final}
        & \scalarinput{combined_dollar_basic_credit_final}
        & \scalarinput{combined_dollar_premium_credit_final} \\
      \multicolumn{2}{>{\raggedright}p{1in}}{Homogeneous}
        & \scalarinput{combined_dollar_cash_naive}
        & \scalarinput{combined_dollar_regulated_debit_naive}
        & \scalarinput{combined_dollar_unregulated_debit_naive}
        & \scalarinput{combined_dollar_basic_credit_naive}
        & \scalarinput{combined_dollar_premium_credit_naive} \\
      \hline
      \multicolumn{2}{>{\raggedright}p{0.9in}}{Difference}
        & \scalarinput{combined_dollar_cash_overstate}
        & \scalarinput{combined_dollar_regulated_debit_overstate}
        & \scalarinput{combined_dollar_unregulated_debit_overstate}
        & \scalarinput{combined_dollar_basic_credit_overstate}
        & \scalarinput{combined_dollar_premium_credit_overstate} \\
       & Consumer Sorting
        & \scalarinput{combined_dollar_cash_sorting}
        & \scalarinput{combined_dollar_regulated_debit_sorting}
        & \scalarinput{combined_dollar_unregulated_debit_sorting}
        & \scalarinput{combined_dollar_basic_credit_sorting}
        & \scalarinput{combined_dollar_premium_credit_sorting} \\
       & Merchant Fee Heterogeneity
        & \scalarinput{combined_dollar_cash_firm}
        & \scalarinput{combined_dollar_regulated_debit_firm}
        & \scalarinput{combined_dollar_unregulated_debit_firm}
        & \scalarinput{combined_dollar_basic_credit_firm}
        & \scalarinput{combined_dollar_premium_credit_firm} \\
      \hline
    \end{tabular}
  \end{subtable}

  \captionsetup{font=small}
  \caption*{\emph{Notes:} These tables summarize the consumer-welfare effects of setting all interchange
  fees to zero for consumers with different payment preferences. Panel \ref{tab:redist-bps} shows the
  effects measured as a percentage of each group's baseline consumption. ``Homogeneous merchants'' is the
  effect from a model in which consumers do not sort across merchants, so the composition of payments is
  identical across merchants. ``Full model'' is the effect from our framework that incorporates both
  heterogeneity in the composition and the costs of payments across merchants. ``Consumer Sorting'' captures
  the effects from accounting for variation in the composition of payments across merchants. ``Merchant Fee
  Heterogeneity captures the combined effects of fee variation across sectors and by firm size. Panel
  \ref{tab:redist-dollars} shows the effects in dollar terms, where we normalize total consumer purchases to
  \$12~trillion \citet{Nilson2023a}.}
\end{table}
\par\end{center}

\begin{center}
\pagebreak\begin{table}[H]
  \centering
  \caption{Effects of the Durbin Amendment}
  \label{tab:durbin}
\begin{tabular}{lccccc}
\toprule
\multirow{2}{*}{Change in Consumption} & \multicolumn{5}{c}{Consumer Group} \\
\cmidrule(lr){2-6}
& Cash & \makecell{Regulated \\ Debit} & \makecell{Exempt \\ Debit}  & \makecell{Basic \\ Credit}  & \makecell{Premium \\ Credit} \\
\midrule
Basis Points & \scalarinput{durbin_cash} & \scalarinput{durbin_regulated_debit} & \scalarinput{durbin_unregulated_debit} & \scalarinput{durbin_basic_credit} & \scalarinput{durbin_premium_credit} \\
Dollars (\$Bn) & \scalarinput{durbin_cash_dollar} & \scalarinput{durbin_regulated_debit_dollar} & \scalarinput{durbin_unregulated_debit_dollar} & \scalarinput{durbin_basic_credit_dollar} & \scalarinput{durbin_premium_credit_dollar} \\
\bottomrule
\end{tabular}
  \caption*{\emph{Notes:} Table \ref{tab:durbin} summarize the consumer-welfare effects of capping the interchange fees 
  on regulated debit cards as per the Durbin Amendment.}
\end{table}

\par\end{center}

\begin{center}
\pagebreak\begin{table}[H]
  \centering
  \caption{The Rise of Premium Credit Cards}
  \label{tab:premiumization}
\begin{tabular}{lccccc}
\toprule
\multirow{2}{*}{Change in Consumption} & \multicolumn{5}{c}{Consumer Group} \\
\cmidrule(lr){2-6}
& Cash & \makecell{Regulated \\ Debit} & \makecell{Exempt \\ Debit}  & \makecell{Basic \\ Credit}  & \makecell{Premium \\ Credit} \\
\midrule
Basis Points of Consumption & \scalarinput{premiumization_cash} & \scalarinput{premiumization_regulated_debit} & \scalarinput{premiumization_unregulated_debit} & \scalarinput{premiumization_basic_credit} & \scalarinput{premiumization_premium_credit} \\
Dollars (billions) & \scalarinput{premiumization_cash_dollar}& \scalarinput{premiumization_regulated_debit_dollar} & \scalarinput{premiumization_unregulated_debit_dollar} & \scalarinput{premiumization_basic_credit_dollar} & \scalarinput{premiumization_premium_credit_dollar} \\
\bottomrule
\end{tabular}
  \caption*{\emph{Notes:} Table \ref{tab:premiumization} summarize the consumer-welfare effects of the rise in premium credit cards.}
\end{table}

\par\end{center}

\newpage{}

\appendix

\section{Additional Tables}

\renewcommand{\thetable}{A\arabic{table}} 
\setcounter{table}{0}        
\renewcommand{\thefigure}{A\arabic{figure}} 
\setcounter{figure}{0}            
\begin{center}
\begin{table}[H]
\caption{Effects of the Durbin Amendment in states with no-surcharge rules\label{tab:durbin-surcharging}}

\begin{centering}
{\tiny\input{../output/tables/durbin_iv_nosurcharge}}
\par\end{centering}
\caption*{\emph{Notes: }This table shows the results of regressions
studying the cross-sectional effects of interchange fees on merchant
sales in the $10$ states that had no-surcharge rules at the time
of the Durbin Amendment. The selected states are California, Colorado,
Connecticut,Florida, Kansas, Massachusetts, Maine, New York, Oklahoma,
and Texas.}
\end{table}
\par\end{center}
\end{document}
